\documentclass[a4paper,twoside,12pt]{report}
% Richard Klein (2020,2021)

% Include Packages
%\usepackage[a4paper,inner=3.5cm,outer=2.5cm,top=2.5cm,bottom=2.5cm]{geometry}  % Set page margins
\usepackage{tikz}
\usepackage[font=small]{caption}
\usepackage{subcaption}
\usepackage{fullpage}
\usepackage{float}                  % Allows 'Here and Only Here' [H] for Floats
\usepackage{url}                    % \url{} command
\usepackage{charter}                  % Set font to Times
\usepackage{graphicx}               % \includegraphics
\usepackage{amsmath}
\usepackage{amssymb}
\usepackage{amsthm}
\usepackage{booktabs}
\usepackage{parskip}
\usepackage{ragged2e}
\usepackage[all]{nowidow}
\setnoclub[2]
\setnowidow[2]

% Referencing
% Provides \Vref and \vref to indicate where a reference is.
\usepackage{varioref} 
% Hyperlinks references
\usepackage[bookmarks=true,bookmarksopen=true]{hyperref} 
% Provides \Cref, \cref, \Vref, \vref to include the type of reference: fig/eqn/tbl
\usepackage{cleveref} 
% Setup Hyperref
\hypersetup{
  colorlinks   = true,              %Colours links instead of ugly boxes
  urlcolor     = blue,              %Colour for external hyperlinks
  linkcolor    = blue,              %Colour of internal links
  citecolor    = blue                %Colour of citations
}
% Names for Clever Ref
\crefname{table}{table}{tables}
\Crefname{table}{Table}{Tables}
\crefname{figure}{figure}{figures}
\Crefname{figure}{Figure}{Figures}
\crefname{equation}{equation}{equations}
\Crefname{equation}{Equation}{Equations}

% Wits Citation Style
\usepackage{utils/natbib} % Force natbib.sty to put citation labels in the reference list
\makeatletter
\renewcommand\NAT@biblabel[1]{\def\citeauthoryear##1##2{##1 ##2}[#1]\hfill}
\renewcommand\NAT@bibsetup[1]{%
  \setlength{\itemsep}{\bibsep}\setlength{\parsep}{\z@}}
\def\@lbibitem[#1]#2{%
  \if\relax\@extra@b@citeb\relax\else
    \@ifundefined{br@#2\@extra@b@citeb}{}{%
     \@namedef{br@#2}{\@nameuse{br@#2\@extra@b@citeb}}}\fi
   \@ifundefined{b@#2\@extra@b@citeb}{\def\NAT@num{}}{\NAT@parse{#2}}%
   \item[\hfil\hyper@natanchorstart{#2\@extra@b@citeb}\@biblabel{#1}%
    \hyper@natanchorend]%
    \NAT@ifcmd#1(@)(@)\@nil{#2}}
\makeatother


\bibliographystyle{utils/named-wits}
\bibpunct{[}{]}{;}{a}{}{}  % to get correct punctuation for bibliography
\setlength{\skip\footins}{1.5cm}
\newcommand{\citets}[1]{\citeauthor{#1}'s \citeyearpar{#1}}
\renewcommand\bibname{References}  

\pagestyle{headings}

\pagestyle{plain}
\pagenumbering{roman}

\renewenvironment{abstract}{\ \vfill\begin{center}\textbf{Abstract}\end{center}\addcontentsline{toc}{section}{Abstract}}{\vfill\vfill\newpage}
\newenvironment{declaration}{\ \vfill\begin{center}\textbf{Declaration}\end{center}\addcontentsline{toc}{section}{Declaration}}{\vfill\vfill\newpage}
\newenvironment{acknowledgements}{\ \vfill\begin{center}\textbf{Acknowledgements}\end{center}\addcontentsline{toc}{section}{Acknowledgements}}{\vfill\vfill\newpage}

\begin{document}
\RaggedRight % Enforce Wits readability rule (avoid full justification)
\pagenumbering{roman} % Roman numerals for preliminaries
\thispagestyle{empty}

% -------------------------
% TITLE PAGE
% -------------------------
\begin{center}
  % \vspace*{2cm}
  {
  \huge \bf \textsc{CRISP-NAM for Survival Analysis under Competing Risks using Risk Set Sampling}\\
  \vspace{1.5cm}
  \large \textbf{Michael Falk}\\
  3052851\\[1.5cm]
  \includegraphics[width=2cm]{utils/images/wits.jpg} % Adjust path as needed
  \vspace{1.5cm}\\
  \large A research report submitted to the Faculty of Science, University of the Witwatersrand, Johannesburg, in partial fulfilment of the requirements for the degree of Master of Science.\\[1.5cm]
  Supervised by Dr. Justine Nasejje\\[5pt]
  Co-Supervised by Dr. Steven James\\[1.5cm]
  Johannesburg, \today\\
  \vspace{0.5cm}
  \small{Ethics Clearance Number: WCSAM-2026-25}\\
  }
\end{center}
\vfill
\newpage

% -------------------------
% DECLARATION (Wits Standard Wording)
% -------------------------
\begin{declaration}
I declare that this research report is my own, unaided work. It is being submitted for the Degree of Master of Science at the University of the Witwatersrand, Johannesburg. It has not been submitted before for any degree or examination at any other University.

\vspace{2cm}
\noindent\rule{6cm}{0.4pt}\\
(Signature of candidate)\\

\vspace{1cm}
\noindent \rule{1cm}{0.4pt} day of \rule{3cm}{0.4pt} 20\rule{1cm}{0.4pt} at \rule{4cm}{0.4pt}
\end{declaration}

% -------------------------
% ABSTRACT
% -------------------------
\begin{abstract}
% Wits Rule: Max 150 words for MSc Research Report, Max 350 for PhD.
Abstract text goes here.
\end{abstract}

% -------------------------
% ACKNOWLEDGEMENTS
% -------------------------
\begin{acknowledgements}
I would like to acknowledge Dr. Justine Nasejje for assisting and guiding me in the creation of this study, as well as all the lecturers who imparted their knowledge and time into my understanding of the work.
\end{acknowledgements}

% -------------------------
% LISTS (TOC, Figures, Tables, Nomenclature)
% -------------------------
\tableofcontents
\newpage

% \phantomsection
% \addcontentsline{toc}{chapter}{List of Figures}
% \listoffigures
% \newpage

% \phantomsection
% \addcontentsline{toc}{chapter}{List of Tables}
% \listoftables
% \newpage

% ======================================================================
% BODY CHAPTERS
% ======================================================================
\clearpage
\pagenumbering{arabic} % Switch to Arabic numerals for main text

\chapter{Introduction}
\label{ch:introduction}

\section{Background and Motivation}

The modern world is currently navigating a highly data driven era, where the volume, velocity, and variety of information generated across multiple domains continue to grow at an unprecedented, exponential rate \cite{Laney2001}. From global financial transactions and sensor networks to electronic health records and longitudinal patient tracking, there is now a digital footprint for almost every human interaction. While this vast availability of data has paved the way for highly sophisticated, deep learning analytical models \cite{LeCun2015} , it has simultaneously introduced massive computational challenges. Training algorithms to process and extract meaningful patterns from these sprawling datasets requires substantial hardware power, highlighting the urgent need for scalable processing techniques.

The computational burden is aslo equaly apparent in the medical industry, particularly within the domain of survival analysis. Survival analysis evaluates the time until a specific event of interest occurs, such as the onset of a disease or patient mortality \cite{Cox1972} . A fundamental component of this analysis is the concept of the \textit{risk set} this is a pool of subjects who are still event free and under observation at any given point in time. As the size of a clinical cohort grows, the mathematical risk set grows proportionally. In modern deep survival models, comparing an individual experiencing an event against the entire available risk set scales quadratically ($O(N^2)$), turning the evaluation of large datasets into a severe computational bottleneck.

This challenge is further compounded in real-world clinical settings by the presence of \textit{competing risks} \cite{latouche2013competing}. In a competing risks scenario, a patient may experience an alternative, mutually exclusive event that permanently precludes the primary event of interest from ever occurring (for example, a patient passing away from a non-cardiovascular cause before they can develop coronary heart disease) \cite{andersen2012competing}. To accurately model the cumulative incidence of the primary event without introducing statistical bias, models commonly rely on the Fine-Gray subdistribution hazard framework \cite{FineGray1999}. However, the Fine-Gray formulation mathematically forces subjects who experience a competing event to remain in the risk set indefinitely, adjusted via censoring weights. As a result, the risk set does not organically shrink over time as it would in a standard Cox model; instead, it remains dense and massive, making full-risk training computationally paralyzing for deep neural networks \cite{Katzman2018, Lee2018}.

As deep survival architectures become more complex to capture non-linear clinical relationships \cite{kvamme2019}, their training durations have increased drastically. This extended training time is highly detrimental to the machine learning lifecycle. Model development is an inherently iterative process: an architecture must be trained before it can be evaluated, its predictions tested, and its hyperparameters tuned. If a single training run takes an excessive amount of time, or if a model trains incorrectly and must be restarted, the entire research, evaluation, and deployment pipeline is delayed. Speeding up the training phase directly accelerates the testing and validation cycles.

To resolve this, this thesis investigates the application of risk set sampling techniques such as Nested Case Control, Static Case-Cohort and the study's own dynamic sampling strategy to deep competing-risk models. By shifting away from training on the exhaustive full-cohort risk set and instead training on a strategically sampled, information-dense approximation of the cohort, it is possible to drastically accelerate model training \cite{Breslow1974, Prentice1986, Borgan2000}. Furthermore, as the world faces increasing constraints on global computing resources, memory bandwidth, and CPU/RAM availability in 2026, the ability to train robust deep learning models while dramatically reducing memory footprints and processing overhead is no longer just a convenience it is an operational necessity.

This study proposes that through the use of advanced, dynamic sampling techniques, deep survival models can maintain their statistical validity and predictive parity with full-risk models, all while achieving vital computational efficiency. By establishing a framework that samples intelligently rather than exhaustively, this research aims to streamline the evaluation of competing risks in large-scale medical data.

\section{Research Questions and Objectives}

\subsection{Research Questions}
This study investigates the effects of risk-set sampling on competing risks for right-censored data, along with a custom dynamic sampling technique. The key research questions guiding this thesis are:
\begin{itemize}
    \item How do case-cohort and nested case-control sampling designs affect discrimination, calibration, and explained variation in Fine-Gray models for competing risks compared to full risk analyses?
    \item How does the proposed dynamic sampling strategy, which adapts dynamically, affect discrimination, calibration, and explained variation in Fine-Gray models for competing risks compared to full risk analyses as well as the other techniques (NCC \& SCC)?
    \item What network and training design choices (e.g., additive interpretable modules, weight stabilization, censoring corrections, mini-batch risk-set sampling) ensure stable and consistent training on sampled data?
    \item Between case-cohort, nested case-control, and the study's own dynamic sampling strategy designs, which provides superior performance in terms of discrimination (C-index), Brier score, Integrated Brier score, robustness to censoring, and computational efficiency (CPU Usage, Memory Usage, and Training Durations)?
    \item How do different sampling strategies affect CPU usage, memory requirements, and training duration for deep Fine-Gray models?
\end{itemize}

\subsection{Research Objectives}
To address the research questions outlined above, this study pursues the following primary objectives:
\begin{itemize}
    \item To implement interpretable deep Fine-Gray models (CRISP-NAM) with stabilized inverse probability weighting censored (IPWC) for sampled risk sets using case-cohort, nested case-control, and the proposed dynamic sampling strategy.
    \item To evaluate and compare the predictive performance (discrimination, calibration, and explained variation) of full-cohort, case-cohort, nested case-control, and dynamic sampling designs using right-censored survival data.
    \item To evaluate the proposed dynamic sampling strategy in terms of predictive performance (discrimination, calibration) using C-index and Brier score against the full risk-set model as a statistical baseline, as well as robustness and computational efficiency (CPU usage, memory footprint, and training time), compared to case-cohort and nested case-control methods.
    \item To analyze the impact of different sampling strategies on computational efficiency, including CPU usage, memory requirements, and training duration across model configurations.
    \item To analyze the impact of different sampling strategies on the C-index, Brier Score, and Integrated Brier Score across model configurations.
    \item To provide practical guidelines for network design and training strategies that maintain interpretability, calibration, stability, and computational efficiency when using sampled risk sets.
    \item To produce reproducible code and benchmark results to support the adoption of sampling-aware deep survival analysis methods.
\end{itemize}

\section{Thesis Organization}
This thesis is structured to systematically introduce, develop, and evaluate the proposed sampling methodologies within the context of deep survival analysis. 
\begin{itemize}
    \item \textbf{Chapter 2 (Background Theory):} Explores the theoretical foundations of competing risks, the Fine-Gray subdistribution hazard, traditional risk-set sampling strategies, and modern deep survival architectures.
    \item \textbf{Chapter 3 (Methodology):} Details the preprocessing pipelines, the mathematical formulation of the novel dynamic entropy sampling strategy, the CRISPNAM-FG network architecture, and the quantitative evaluation metrics.
    \item \textbf{Chapter 4 (Results and Discussion):} Presents the empirical findings across both synthetic and real-world datasets, highlighting the trade-offs between predictive discrimination, probability calibration, and computational efficiency. 
    \item \textbf{Chapter 5 (Conclusion and Future Work):} Summarizes the major contributions of the study, addresses limitations, and outlines potential avenues for future research in scalable survival analysis.
\end{itemize}
% --------------------------------------------------------------
\chapter{Background Theory}
\label{ch:background_theory}

This chapter establishes the theoretical foundations required to model time to event data in the presence of multiple, mutually exclusive outcomes. It first reviews the mathematical challenges introduced by competing risks and evaluates classical statistical frameworks, including cause-specific and subdistribution hazards. The chapter then explores recent advancements in deep survival analysis, detailing how neural architectures can be adapted for competing events. Finally, it introduces the CRISPNAM framework and the risk set sampling strategies necessary for computationally efficient model optimization.

\section{Competing Risks in Survival Analysis}
While traditional survival analysis focuses on predicting risk for a single event of interest, many real-world problems such as death or getting infected, particularly in the healthcare domain, involve more than one competing event of interest. For example, patients who are at risk of death due to a specific disease may pass away following other unrelated complications, which acts as a competing risk that must be taken into account when predicting the primary risk. Conventional survival analysis approaches, including the standard Kaplan-Meier estimator and the Cox proportional hazards model \cite{Cox1972}, are inappropriate for competing risk scenarios because they erroneously handle alternative events as standard censored observations. This fundamental misclassification artificially alters the risk pool, resulting in inflated estimates of event probabilities and absolute risk. 

\section{Past Deep Neural Network Models}
There have been a steady growth in the advances in the area of deep learning have demonstrated improved predictive performance for survival analysis by replacing rigid parametric assumptions with flexible neural architectures capable of capturing complex, non-linear patterns in the data \cite{LeCun2015}. A foundational model in this domain is DeepSurv \cite{Katzman2018}, which successfully extended the traditional Cox proportional hazards framework by replacing the standard linear predictor with a deep neural network. The network learns the relationship between a patient's covariates and their risk by optimizing the standard Cox partial-likelihood objective \cite{Katzman2018}. Building upon this neural-Cox integration, Kvamme et al. \cite{kvamme2019} further advanced continuous-time survival predictions by developing neural networks that relax the strict proportional hazards assumption, enabling the model to explicitly learn complex and time-varying feature effects. However, while highly effective for standard survival settings, models like DeepSurv and the architectures proposed by Kvamme et al. are primarily designed for single-event survival and do not natively accommodate the mathematical complexities of competing risks.

To address multiple event types, methodologies such as DeepHit \cite{Lee2018} have been developed. DeepHit presents a joint modeling framework that employs a shared representation network coupled with cause-specific sub-networks to model the joint distribution of event time and type. The model optimization intelligently combines negative log-likelihood with a ranking loss to ensure concordance between predicted and observed outcomes \cite{Lee2018}. While DeepHit achieves strong predictive gains in competing risk scenarios, these gains come at the cost of interpretability, resulting in a black-box model that is difficult for clinicians to trust and validate in clinical practice.

\section{Cause Specific Hazard}
One of the most classical methodologies for estimation in competing risks survival analysis is the Cause Specific Hazard method, originally proposed by Prentice et al. \cite{prentice1978analysis}. Building directly upon the foundational Cox proportional hazards framework \cite{Cox1972}, this method models competing events separately by treating each outcome as a distinct hazard function. Under this formulation, the model evaluates the instantaneous rate of the primary event, conditional on the patient having survived completely event-free up to that time. 

For each cause $k$, the cause-specific hazard $\lambda_k(t|x)$ is mathematically defined as the instantaneous risk of experiencing event type $k$ at time $t$, given that the individual has survived without any event up until time $t$:
$$\lambda_k(t|x) = \lim_{\Delta t \to 0} \frac{P(t \le T < t+\Delta t, E = k \mid T \ge t, x)}{\Delta t}$$

Traditionally, the cause-specific hazard function is used in conjunction with the Cumulative Incidence Function (CIF) to calculate the risk score and determine the probability of event $k$ occurring over time. The CIF represents the probability that event type $k$ occurs by time $t$, accounting for the presence of competing risks. It is computed by integrating the cause-specific hazard weighted by the overall event-free survival function ($S(s|x)$):
$$F_k(t|x) = \int_0^t \lambda_k(s|x) S(s|x) ds$$

However, these classical approaches including standard Cox proportional hazards models heavily rely on risk-set definitions that erroneously treat alternative events as censored observations \cite{Cox1972}. This is problematic because it equates competing risks with standard right censoring. Right censoring is the most common form of censoring found in medical and reliability survival analysis, that is when the event of interest has not been observed for an individual during the study period \cite{kalbfleisch2002statistical}. In other words, we know the patient survived at least up to a certain point in time, but their exact outcome or the exact time of an event after that remains unknown. 

Consequently, subjects experiencing competing events are permanently removed from the risk set and treated mathematically as these right-censored observations \cite{latouche2013competing}. While computationally straightforward, cause-specific approaches impose strict independence assumptions. By treating competing events merely as censoring, the method fails to capture the true risk dynamics, ultimately leading to inflated and biased estimates of absolute event probabilities when competing events are dependent or when precise absolute risk predictions are required \cite{pintilie2006competing, andersen2012competing}.


\section{Fine-Gray Subdistribution Hazard}

To properly estimate absolute risk without the biases inherent to the cause-specific model, the Fine-Gray subdistribution method was introduced \cite{FineGray1999}. This method directly models the subdistribution hazard, maintaining subjects who experience competing events in the risk set. 

As an example of this contrast, in standard Cox models, competing events are treated as censored and the individuals are removed from the risk set. In the Fine-Gray model, individuals who have experienced competing events remain in the risk set but with a modified contribution. The subdistribution risk pool includes both subjects who are event-free at time $t$ and those who experienced a competing event before time $t$. Because these patients can no longer experience the primary event of interest, their influence in the risk pool is adjusted over time through Inverse Probability of Censoring Weights (IPCW), utilizing the Kaplan-Meier estimate of the censoring distribution.

The subdistribution hazard is denoted in two terms:
$$\tilde{\lambda}_k(t|X)=\tilde{\lambda}_{k,0}(t)\exp(\beta^\top X)$$
where $\tilde{\lambda}_{k,0}(t)$ is the baseline subdistribution hazard, which is used to capture the time dependence of the event $k$, and $\exp(\beta^\top X)$ is the relative risk of the individual based on the covariates $X$.

By keeping these patients in the risk pool, the Fine-Gray model preserves the probability structure induced by competing outcomes, enabling the direct modeling of the Cumulative Incidence Function (CIF). As we can see again, we have a CIF with the Fine-Gray subdistribution hazard that accomplishes the same task as with the cause-specific hazard function: to calculate the probability that event $k$ occurs by time $t$ while accommodating for the competing risk. It is mathematically expressed as:
$$F_k(t|X) = P(T \le t, \Delta=k|X)$$

A key feature of the Fine-Gray approach is the use of a modified risk set, denoted by an indicator such as $Y_i^*(t)$. This indicator differs from standard survival analysis because it retains individuals who have experienced competing events, rather than removing them entirely. The weighted risk contribution describes how the Fine-Gray model applies appropriate weighting within the risk set to ensure correct representation in the estimation procedure, accounting for individuals who are censored or have experienced competing events. 

The individuals are then put in the modified risk set that is captured through $Y_i^*(t)$. It is utilized to compute $\bar{x}(t; \beta)$, which represents the risk-set weighted average of the covariates at time $t$. It is calculated over the modified risk set using IPCW weights $w_i(t)$ proportional to the relative risk $\exp(\beta^\top X)$, and it appears in the score function used to estimate the regression parameters:
$$\bar{x}(t; \beta)=\frac{\sum_i Y_i^*(t) w_i(t) X_i \exp(\beta^\top X_i)}{\sum_i Y_i^*(t) w_i(t) \exp(\beta^\top X_i)}$$


\section{CRISPNAM}

While traditional deep learning approaches offer high predictive capacity, their inherent "black-box" nature often limits clinical interpretability. A foundational work shaping the theoretical framework of this thesis is the approach by Ramachandram et al. \cite{InterpretableFineGray2025}, which introduced an interpretable Fine-Gray deep survival model designed specifically for competing risks. Building upon this vital need for transparency in deep survival analysis, the CRISPNAM architecture leverages a Neural Additive Model (NAM) framework to achieve both high flexibility and strict interpretability.

Central to the CRISPNAM architecture is the use of the FeatureNet. Unlike standard fully connected neural networks that process all inputs simultaneously through hidden layers, the additive framework assigns a dedicated, independent multi-layer perceptron a FeatureNet to each individual covariate. Let $X = (x_1, x_2, \dots, x_p)$ represent the set of $p$ input features. Each feature $x_j$ is passed exclusively through its own FeatureNet, denoted as $f_j(x_j)$, which learns a specific non-linear shape function for that variable alone. 

The additive framework mathematically combines these independent learned functions by summing them together to construct the overall risk predictor $f(X)$:
$$f(X) = \sum_{j=1}^p f_j(x_j)$$

This formulation elegantly integrates with the Fine-Gray method. By replacing the traditional linear predictor $\beta^\top X$ with the learned additive neural predictor $f(X)$, the modified subdistribution hazard becomes:
$$\tilde{\lambda}_k(t|X) = \tilde{\lambda}_{k,0}(t)\exp\left(\sum_{j=1}^p f_j(x_j)\right)$$

Because the model computes the hazard additively, the exact contribution of any single feature $x_j$ to the overall relative risk can be isolated and visualized. This architecture allows the model to capture complex, non-linear biostatistical relationships while retaining the exact interpretability of traditional regression models.


\section{Risk Set Sampling Strategies}

While the Fine-Gray modified risk set mathematically solves the statistical bias of competing risks, retaining individuals indefinitely creates a severe computational burden. Evaluating the full risk set at every event time scales quadratically, becoming a severe bottleneck for large datasets \cite{Wang2020}. To mitigate this, traditional biostatistics employ risk set sampling designs to computationally approximate the full cohort. 

One such method is Nested Case-Control (NCC) sampling, a time-matched design where, for each observed case $i$ at a given event time $t_i$, a fixed quota of $m$ controls is sampled uniformly without replacement from the valid risk set active at that specific moment, denoted as $R(t_i)$. Because controls are matched sequentially to specific events, it creates a much smaller, localized approximation of the risk set, $\tilde{R}(t_i)$. The probability of a control $j$ being included fluctuates depending entirely on the varying size of the full risk set at that time, $n(t_i)$. This inclusion probability is mathematically expressed as:
$$P(j \in \tilde{R}(t_i) \mid j \in R(t_i), j \neq i) = \frac{m}{n(t_i) - 1}$$

Alternatively, the Static Case-Cohort (SCC) design operates independently of event times by selecting a global subcohort $C$ of patients randomly at baseline (time zero) using a fixed sampling fraction $\alpha$ \cite{Prentice1986}. This static subcohort is then reused throughout the entirety of the training process. At any given event time $t_i$, the evaluated risk set consists of the current case $i$ along with any members of the baseline subcohort who are still actively at risk:
$$\tilde{R}(t_i) = \{i\} \cup (C \cap R(t_i))$$
Under this design, the probability of any individual being initially selected into the subcohort is strictly determined by the fixed baseline fraction:
$$P(j \in C) = \alpha$$

\section{Inverse Probability Weighting (IPW) and Stabilization}

While sub-sampling architectures such as nested case-control and case-cohort designs significantly alleviate computational overhead, they fundamentally alter the underlying data distribution by over representing cases relative to controls \cite{Prentice1986, Borgan2000}. A naive analysis of this skewed subset would yield biased predictions. To rectify this sampling bias, Inverse Probability Weighting (IPW) is applied \cite{HorvitzThompson1952}. 

This technique scales the contribution of each selected observation by the reciprocal of its sampling probability, allowing the subset to accurately mirror the full cohort. If $\pi_i$ represents the probability that individual $i$ is selected into the sample, the standard IPW is expressed mathematically as:
$$w_i = \frac{1}{\pi_i}$$
Consequently, individuals who possess a lower likelihood of being sampled are assigned proportionally larger weights. This mechanism ensures that the sampled cohort yields consistent and unbiased estimates of the true population parameters \cite{Prentice1986, Borgan2000}.

However, a critical vulnerability of standard IPW arises when inclusion probabilities ($\pi_i$) become exceedingly small, generating disproportionately large weights. In the context of machine learning and gradient-based optimization, this extreme variability introduces significant numerical instability and severely inflates the variance of the estimator. To counter this, statistical and causal inference frameworks frequently employ weight stabilization \cite{Robins2000, Hernan2000}. By scaling the baseline weight by the overall marginal probability of selection into the sample, $P(S=1)$, the stabilized weight is defined as:
$$w_i^{\text{stab}} = \frac{P(S=1)}{\pi_i}$$
This mathematical normalization centers the weights and drastically curtails empirical variance. It preserves the bias-correcting properties of standard IPW while ensuring smoother convergence and computational stability during neural network training.

Beyond intentional sampling disparities, survival datasets are routinely subject to right-censoring, where the exact time of the event is not recorded for all individuals. To handle subjects whose true event times are partially unobserved, the framework incorporates Inverse Probability of Censoring Weighting (IPCW) \cite{Graf1999}. IPCW dynamically adjusts the contribution of subjects based on the inverse of their probability of remaining uncensored up to a specific time, a probability that is typically derived from a censoring survival estimator, denoted as $\hat{G}(t|X)$. 

By integrating both stabilization and IPCW into the objective function, the model guarantees that inferences remain statistically valid and reflective of the true population dynamics, even when trained on partially observed and heavily sub-sampled data.


\chapter{Methodology}
\label{ch:methodology}

This chapter details the experimental design, data processing pipelines, and computational frameworks used to evaluate the proposed CRISPNAM architecture along with the  incoperation of the smapling techniques, that also include our custome sampling strategy. It begins by outlining the acquisition and preprocessing of both real-world clinical datasets and synthetically generated cohorts designed to test specific mathematical assumptions. The chapter then details the implementation of the deep survival models, including the network architecture, loss function formulations, and the risk set sampling techniques. Finally, it establishes the evaluation metrics and cross-validation strategies used to rigorously benchmark predictive performance .

\section{Data Preprocessing}

To evaluate the robustness, accuracy, and interpretability of the competing risks models, experiments were conducted across three distinct data environments: a fully controlled synthetic dataset, the clinical Primary Biliary Cholangitis (PBC) dataset, and a semi-synthetically modified version of the Framingham Heart Study dataset.

\subsection{Synthetic Data Generation}
A fully synthetic dataset was computationally generated to provide a controlled environment where the true underlying data generating distributions, baseline hazards, and covariate coefficients were known definitively. This ground-truth knowledge allows for exact empirical validation of the model's ability to recover true risk estimators without the confounding variables inherent to real-world medical data.

A cohort of $N = 15,000$ independent samples was generated. To ensure full experimental reproducibility, a fixed global random state (seed = 42) was utilized throughout the generation process. The feature space was constructed using 12 covariates: six continuous variables drawn from a standard normal distribution $\mathcal{N}(0, 1)$, and six binary variables sampled from a binomial distribution $B(1, 0.5)$. These were concatenated to form the complete covariate matrix $X$.

Two competing events were simulated. For each event type $k \in \{1, 2\}$, a specific vector of true coefficients $\beta_k$ was predefined. To test the model's capacity for intrinsic feature selection and to simulate uninformative clinical noise, several covariates in both $\beta_1$ and $\beta_2$ were deliberately assigned a true coefficient of zero. The linear predictors for each individual were calculated as $\eta_{ik} = X_i^\top \beta_k$.

The cause-specific hazards were then defined exponentially as:
$$\lambda_k(t|X_i) = \lambda_{0,k} \exp(\eta_{ik})$$
where the baseline hazards were fixed at $\lambda_{0,1} = 0.01$ and $\lambda_{0,2} = 0.05$. 

The total instantaneous hazard for any event occurring was computed as the sum of the cause-specific hazards ($\lambda_{\text{total}} = \lambda_1 + \lambda_2$). The exact time to event for each subject was simulated via inverse transform sampling, drawing from an exponential distribution parameterized by the total hazard:
$$T_{\text{event}} = \frac{-\ln(U)}{\lambda_{\text{total}}}$$
where $U \sim \text{Uniform}(0,1)$. If an event occurred, the specific event type (1 or 2) was assigned probabilistically using a binomial trial, where the probability of experiencing Event 1 was strictly proportional to its relative contribution to the total hazard ($\lambda_1 / \lambda_{\text{total}}$).

Finally, right-censoring was introduced independently to mimic study dropout or administrative study end. Censoring times, $T_{\text{censor}}$, were drawn from an exponential distribution with a predefined censoring rate of $0.01$. The final observed survival time was derived as the minimum of the simulated event time and the censoring time ($\tilde{T} = \min(T_{\text{event}}, T_{\text{censor}})$), and the event indicator was recorded as 0 if the individual was censored, or assigned the respective event type otherwise.

\subsection{Primary Biliary Cholangitis (PBC) Dataset}
The Primary Biliary Cholangitis (PBC) dataset, sourced from the Mayo Clinic trial, provides a classical real-world baseline for evaluating competing risks. Data preprocessing for this cohort initially required standardizing the temporal scale; observation times were converted from years into days to provide a higher-resolution continuous time axis for the neural network. 

The clinical outcomes were mapped into a defined competing risks state space: State 0 denoted right-censored patients, State 1 denoted death (the primary event of interest), and State 2 denoted liver transplantation (the competing event). Addressing missingness in the clinical covariates was crucial for tensor continuity. Missing continuous variables (such as serum cholesterol or prothrombin time) were imputed using the cohort mean. Categorical covariates were imputed using the mode (the most frequent class) prior to being transformed via one-hot encoding. Finally, redundant administrative identifiers and any medical records lacking valid, positive survival times were systematically excluded to establish the final analytical cohort.

\subsection{Framingham Heart Study Dataset}
To evaluate the models on a highly dimensional and well-understood cardiovascular cohort, the standard Framingham Heart Study dataset was adapted. Initial data cleaning involved resolving missing values via median imputation, ensuring that feature distributions were not severely skewed by outliers. 

Because standard Framingham data is often structured for binary classification rather than time-to-event analysis, the original deterministic target column (\texttt{ten\_year\_c\_h\_d}) was explicitly dropped from the feature matrix to eliminate any risk of data leakage during model optimization. As the baseline dataset lacked precise longitudinal tracking for multiple event types, time-to-event data was semi-synthetically generated to classify patient trajectories into three mutually exclusive states.

Event 1 (The Main Event) represented patients who actively developed Coronary Heart Disease (CHD). Their time-to-event was distributed randomly within the 10-year observational window. To induce a mathematically valid competing risk scenario, Event 2 (The Competing Event) was engineered to simulate a scenario where approximately 15\% of the non-CHD patients experienced a terminal occurrence such as a fatal accident or non-cardiovascular mortality prior to the 10-year mark. 

This semi-synthetic modification is vital for assessing the Fine-Gray loss function, as experiencing this alternative Event 2 permanently precludes the patient from ever experiencing the main CHD event, fundamentally shifting the risk set dynamics. Finally, Event 0 (Censoring) captured the remaining baseline patients who either safely survived event-free to the end of the 10 year study or were randomly lost to follow-up.

\section{Risk Set Sampling}

\subsection{Full Risk}
In standard survival models, computing the objective function requires evaluating the full risk set at every unique event time. When analyzing massive clinical datasets, this exhaustive evaluation becomes computationally prohibitive. To alleviate this bottleneck, various sampling architectures are implemented to condense the risk set at any given time step, actively reducing the number of subjects processed while theoretically preserving the integrity of the Fine-Gray loss function. By randomly drawing a subset of controls to accompany the observed case at time $t_j$ from the valid pool of subjects where $T_i \ge t_j$ \cite{Breslow1974}, computational complexity is drastically reduced while maintaining the localized, time-dependent probability structure of the subdistribution hazard.

\subsection{Case Cohort}
The static case-cohort design offers an alternative that operates independently of event times. In this architecture, a fixed subcohort is randomly selected at baseline (time zero) and systematically reused across the entire training lifecycle. While cases are universally retained when they experience the event of interest, the controls are exclusively drawn from this baseline subcohort. Consequently, this introduces uniform, yet unequal, inclusion probabilities where every member of the subcohort shares a constant probability of selection, denoted by a fixed sampling fraction $\pi_i = f$. To rectify the inherent selection bias caused by over representing cases, these fixed probabilities are integrated directly into an inverse probability weighted formulation of the Fine-Gray objective.

\subsection{Nested Case Control}
Conversely, nested case-control sampling employs a time-dependent matching strategy. At each precise event time $t_j$, the observed case is paired with a specific quota of $m$ controls, which are drawn strictly from the active risk set at that exact moment. Because the size of the risk set shrinks over time as patients experience events or are censored, the probability of a control being sampled, denoted as $\pi_{i,j}$, dynamically fluctuates. To ensure the sampled subset unbiasedly mimics the full cohort's structural contribution, these time-varying inclusion probabilities must be explicitly neutralized via inverse probability weighting within the partial likelihood function.

\subsection{Dynamic Sampling Strategy}
Moving beyond classical methodologies, this thesis introduces an adaptive, entropy-guided sampling framework applied directly to the subdistribution risk set. In the context of the Fine-Gray model, the modified risk pool at time $T_i$, denoted $\tilde{R}(T_i)$, mathematically retains individuals who have already experienced competing events, adjusting their influence via Inverse Probability of Censoring Weights (IPCW), $w_j(T_i)$.

This dynamic strategy begins by partitioning the active risk set into distinct groups $\{G_g(T_i)\}_{g=1}^{G_i}$ based on shared covariate profiles, represented by a covariate vector $\boldsymbol{X}_g$. The cumulative effective size of any specific group is calculated by summing its censoring weights:
$$W_g(T_i) = \sum_{j \in G_g(T_i)} w_j(T_i)$$
The empirical proportion of that group relative to the entire modified risk set is then calculated as:
$$p_g(T_i) = \frac{W_g(T_i)}{\sum_h W_h(T_i)}$$

To intelligently target informative yet potentially rare clinical profiles while leveraging the model's current predictive state, an importance score is generated for each group. This score fuses Shannon information gain with the model's current risk assessment:
$$s_g(T_i) = \bigl[-\log p_g(T_i)\bigr]\exp(\hat{\eta}(\boldsymbol{X}_g))$$
where $\hat{\eta}(\boldsymbol{X}_g)$ acts as the current neural risk prediction.

Rather than forcing a rigid sampling quota like the nested case-control design, this approach scales the importance scores against the average group importance to dictate independent inclusion probabilities, bounded by a baseline retention rate $\rho \in (0, 1]$:
$$\pi_g(T_i) = \min \left( 1, \;\rho \cdot \frac{s_g(T_i)}{\frac{1}{G_i}\sum_{h=1}^{G_i} s_h(T_i)} \right)$$
Through this mechanism, highly informative groups ($s_g$ significantly exceeding the mean) are deterministically retained ($\pi_g = 1$), whereas redundant or statistically dense groups are aggressively down-sampled. Because the sampling volume is fluid, the computational batch size organically expands or contracts based on the current informational density of the risk set. 

The resulting sampled collection of groups is denoted as $\tilde{\mathcal{G}}(T_i)$. The contribution of this intelligently sampled subset to the partial likelihood is then estimated using IPW:
$$\widehat{\mathcal{L}}(T_i) = \sum_{g \in \tilde{\mathcal{G}}(T_i)} \frac{W_g(T_i)}{\pi_g(T_i)} \exp\bigl(\eta(\boldsymbol{X}_g)\bigr)$$
Assuming proper specification of inclusion probabilities, this dynamic estimator remains design-consistent. However, because $\pi_g(T_i)$ relies on the active model state $\hat{\eta}(\cdot)$, it functions as an adaptive estimator, intentionally trading slight bias for significant variance reduction and computational efficiency.

\subsection{IPW and Stability}
Fundamentally, all evaluated sampling strategies such as case-cohort, nested case-control, and the dynamic entropy desig are structurally unified through the Inverse Probability Weighting framework. In each design, a selected individual or group contributes to the final Fine-Gray loss inversely proportional to its likelihood of being sampled. 

Let $\pi_i$ represent the generalized inclusion probability for subject $i$ (or $\pi_g(t)$ for a group). The core mechanism of IPW dictates that, under correct probability specification, the weighted estimator seamlessly recovers the full-cohort objective in expectation. This is grounded in Horvitz-Thompson unbiasedness \cite{HorvitzThompson1952}, defined as:
$$\mathbb{E}\!\left[\frac{\mathbf{1}\{i \in S\}}{\pi_i} \, \phi_i \right] = \phi_i$$
where $\phi_i$ represents any individual contribution to the full dataset objective. This unifying mathematical principle ensures that regardless of how the data is subsetted whether statically, temporally, or dynamically the underlying risk estimates remain mathematically stable and reflective of the true clinical population.

\section{Network Architecture (CRISPNAM-FG)}
\subsection{Network architecture}

The modeling framework follows the CRISP-NAM design \cite{InterpretableFineGray2025}, extended here to accommodate mini-batch risk-set sampling under the competing-risks Fine Gray formulation. The architecture is grounded in a Neural Additive Model (NAM) structure, wherein each covariate acts on the sub-distribution hazard independently through its own dedicated sub-network, termed a \textit{FeatureNet}.

Let $\mathbf{x} = (x_1, x_2, \dots, x_p)^\top \in \mathbb{R}^p$ denote the vector of input covariates for a subject. Each scalar feature $x_i$ ($i = 1, \dots, p$) is mapped through a dedicated feedforward sub-network $f_i(x_i): \mathbb{R} \to \mathbb{R}$, yielding an isolated non-linear contribution to the cause-specific log-subdistribution hazard.

\subsubsection{FeatureNet Formulation}
Each FeatureNet $f_i(\cdot)$ is constructed as a deep multilayer perceptron comprising $L$ hidden layers. Setting the initial layer input to the scalar covariate $z_i^{(0)} = x_i$, the transformation proceeds recursively across hidden layers $l = 1, \dots, L$:
\begin{align}
    u_i^{(l)} &= W_i^{(l)} z_i^{(l-1)} + b_i^{(l)}, \\
    z_i^{(l)} &= \text{Dropout}\left( \text{SiLU}\left(u_i^{(l)}\right); p_{\text{drop}} \right), \quad l = 1, \dots, L,
\end{align}
where $W_i^{(l)} \in \mathbb{R}^{d_l \times d_{l-1}}$ and $b_i^{(l)} \in \mathbb{R}^{d_l}$ denote the layer weights and bias vectors, respectively, and $p_{\text{drop}}$ is the dropout retention rate. The sub-networks utilize the Sigmoid Linear Unit ($\text{SiLU}$) activation function, defined as:
\begin{equation}
    \text{SiLU}(u) = u \cdot \sigma(u) = \frac{u}{1 + e^{-u}}.
\end{equation}
The smooth, continuously differentiable, and non-monotonic nature of $\text{SiLU}$ mitigates vanishing gradients across deeper sub-networks while yielding smooth univariate feature effect curves.

The final representation $z_i^{(L)} \in \mathbb{R}^{d_L}$ is projected to a scalar output via a bias-free linear transformation:
\begin{equation}
    \tilde{f}_i(x_i) = W_i^{(L+1)} z_i^{(L)},
\end{equation}
where $W_i^{(L+1)} \in \mathbb{R}^{1 \times d_L}$. To guarantee model identifiability a core requirement of generalized additive models each FeatureNet enforces a zero-mean batch constraint during training:
\begin{equation}
    f_i(x_{b,i}) = \tilde{f}_i(x_{b,i}) - \frac{1}{B} \sum_{j=1}^B \tilde{f}_i(x_{j,i}),
\end{equation}
for sample $b$ within a mini-batch of size $B$. This centering operation ensures that individual shape functions $\mathbb{E}[f_i(x_i)] = 0$ capture purely relative variations, preventing arbitrary vertical translation between subnetworks.

\subsubsection{Global Additive Aggregation}
The overall model aggregates the individual feature-specific transformations additively with a learnable global intercept $\beta_0 \in \mathbb{R}$:
\begin{equation}
    \eta(\mathbf{x}) = \beta_0 + \sum_{i=1}^p f_i(x_i),
\end{equation}
where $\eta(\mathbf{x})$ denotes the predicted log-subdistribution hazard score (relative risk). Because the final projection layers of the FeatureNets omits additive biases, the global baseline risk is entirely parameterized by $\beta_0$.

All linear layer weights are initialized using Xavier uniform initialization 
\begin{equation}
    W \sim \mathcal{U}\left(-\sqrt{\frac{6}{n_{\text{in}} + n_{\text{out}}}}, \sqrt{\frac{6}{n_{\text{in}} + n_{\text{out}}}}\right)
\end{equation}

with biases initialized to zero. During training, the aggregated risk score $\eta(\mathbf{x})$ is optimized directly against the sampling-weighted Fine Gray partial likelihood loss. For post hoc model interpretation, each univariate shape function $f_i(x_i)$ can be visualized over its covariate support domain to inspect linear or non-linear effects on sub-distribution incidence.


\section{Loss Function}

To optimize the neural relative risk predictor under varying computational constraints, the network is trained by minimizing a sampling-aware extension of the Fine-Gray partial likelihood. This objective function is mathematically designed to integrate the additive neural predictor, correct for right-censored observations, and neutralize the selection bias introduced by risk-set sampling architectures. 

The formulation is constructed by first establishing the baseline Fine-Gray objective applied to the full cohort. Let $\mathcal{R}(T_i)$ denote the complete modified risk set evaluated at event time $T_i$ for the target cause $k$. To account for patients who are independently right-censored, Inverse Probability of Censoring Weighting (IPCW) is applied \cite{Graf1999}. If $\hat{G}(t \mid \mathbf{x})$ represents the Kaplan-Meier estimate of the censoring survival distribution, the censoring weight for a subject $j$ at time $T_i$ is computed as:
$$ w_j^{\text{cens}}(T_i) = \frac{1}{\hat{G}(T_i^- \mid \mathbf{x}_j)} $$
Replacing the traditional linear predictor with the additive neural risk score $\eta_k(\mathbf{x})$, the baseline negative log-partial likelihood evaluated over the full risk set is formulated as:
$$ \mathcal{L}_{\text{FG}}(\beta) = -\sum_{i:\Delta_i = k} w_i^{\text{cens}}(T_i) \left[ \eta_k(\mathbf{x}_i) - \log \sum_{j \in \mathcal{R}(T_i)} w_j^{\text{cens}}(T_i) \exp(\eta_k(\mathbf{x}_j)) \right] $$

While statistically sound, optimizing this full baseline objective across all event times is computationally exhaustive. When sub-sampling architectures (such as nested case-control or static case-cohort designs) are employed, the full risk set $\mathcal{R}(T_i)$ is replaced by a much smaller sampled subset, denoted as $\tilde{\mathcal{R}}(T_i)$. However, this subsampling distorts the natural risk distribution. To recover an unbiased representation of the original cohort, Inverse Probability Weighting (IPW) is applied based on the selection probabilities \cite{HorvitzThompson1952, Prentice1986}.

For a given subject $j$ with an inclusion probability of $\pi_j$, the naive IPW weight is $1/\pi_j$. Because extreme fractional probabilities can destabilize neural network gradients (Introducing varince in the weights of the model), a variance reduction technique is applied by scaling the weights by the empirical marginal probability of selection, $P(S=1)$ \cite{Robins2000, Hernan2000}. This yields the stabilized sampling weight:
$$ w_j^{\text{stab}} = \frac{P(S=1)}{\pi_j} $$

The global contribution weight for any subject $j$ evaluated at time $T_i$ is therefore the product of their stabilized sampling weight and their censoring weight:
$$ w_{ij} = w_j^{\text{stab}} \cdot w_j^{\text{cens}}(T_i) $$

By injecting these combined weights and evaluating only over the sampled subset $\tilde{\mathcal{R}}(T_i)$, the final sampling-aware Fine-Gray objective utilized to train the model is derived as:
$$ \mathcal{L}_{\text{FG}}^{\text{IPW}}(\beta) = -\sum_{i:\Delta_i = k} w_{ii} \left[ \eta_k(\mathbf{x}_i) - \log \sum_{j \in \tilde{\mathcal{R}}(T_i)} w_{ij} \exp(\eta_k(\mathbf{x}_j)) \right] $$
where $w_{ii}$ serves as the combined weight of the observed case at its own event time. 

Notably, explicit weight clipping a technique frequently utilized in reinforcement learning and causal inference to bound extreme gradients is intentionally omitted from this framework. Because the sub sampled risk sets are already heavily reduced, artificially truncating valid IPW weights at arbitrary thresholds would discard critical representation information, re-introducing the very statistical bias the IPW formulation is designed to eliminate. 

Throughout the experimental phase of this thesis, the models were aggressively optimized against this sampling-aware $\mathcal{L}_{\text{FG}}^{\text{IPW}}$ loss function. To ensure robust validation and provide a rigorous experimental control, baseline models were trained alongside them using the exact same weighted loss formulation, but evaluated comprehensively over the full, unsampled risk set.


\section{Training Details}

To guarantee rigorous and unbiased evaluation across the varying datasets, the model's predictive capabilities and sampling mechanics were validated using a five-fold cross-validation framework. To address the inherent imbalance often found in survival datasets where the primary event of interest may be significantly overshadowed by censored observations or competing risks a Stratified K-Fold splitting strategy was implemented. By conditioning the fold generation directly on the event indicators, this event-aware stratification ensures that the proportion of competing events, primary events, and right-censored subjects remains strictly consistent across all training and validation splits. This prevents "fold collapse," a phenomenon where a validation set is accidentally starved of minority outcomes, which would severely skew loss gradients and validation metrics. 

Within this cross-validation structure, the experimental pipeline was designed to evaluate both full-batch and mini-batch optimization regimes. This dual regime approach allows for a direct empirical assessment of how disparate sampling architectures interact with varying computational constraints. Full batch training aggregates the partial likelihood across the entirety of the risk set (whether fully enumerated or sampled) before executing a single gradient update. This serves as the classical baseline, mirroring traditional biostatistical optimization. 

Conversely, mini-batch training iterates over strict, localized subsets of the cohort. This mechanism is crucial for scaling deep survival models to massive clinical databases by constraining memory overhead. To seamlessly integrate this into the PyTorch ecosystem, event aware DataLoaders were configured to stream batches while maintaining temporal and categorical stability. However, the convergence dynamics of mini-batching in a competing risks environment are statistically non-trivial. When integrated with aggressive risk-set condensation techniques such as static case-cohort, nested case-control, or the proposed dynamic entropy sampling the neural network is updated using only a fragmented fraction of an already deeply sub-sampled risk pool. 

Consequently, a core objective of this training protocol is to investigate whether this compounded sparsity degrades the stability of the Fine-Gray loss function, or if the Inverse Probability Weighting (IPW) mechanism is robust enough to maintain convergence and predictive accuracy under stochastic mini-batch conditions. This comparative analysis isolates whether sampling alone provides sufficient computational acceleration, or if true scalability requires the hybridization of sampling and mini-batching.

All model configurations were optimized utilizing the Adam optimizer \cite{kingma2014adam}. Adam is highly effective in deep survival contexts as it leverages adaptive learning rates and moving averages of historical gradients to navigate the complex loss topographies generated by IPW scaling. To prevent the FeatureNets from memorizing localized noise a heightened risk when training on deeply sub-sampled risk sets L2 parameter regularization (weight decay) was integrated into the optimization loop. This regularizes the network by penalizing excessively large weights, promoting smoother feature effect curves and ensuring that the interpretable sub-networks generalize effectively to unseen clinical populations.


\section{Model Evaluation and Diagnostics}

This section details the evaluation metrics and diagnostic frameworks employed to validate the predictive performance, calibration, and computational efficiency of the competing risks models. To provide a holistic assessment, the models are evaluated on their discriminative capacity using a time-dependent Concordance Index, their absolute probabilistic accuracy via the Brier and Integrated Brier Scores, and their clinical reliability through Stratified Aalen-Johansen Cumulative Incidence Function (CIF) calibration. Furthermore, computational profiling and synthetic correlation metrics are introduced to explicitly validate the theoretical efficiency and exact feature recovery of the proposed CRISPNAM architecture.

\subsection{Brier Score}
In standard regression tasks, the Mean Squared Error (MSE) evaluates the distance between a prediction and the continuous truth. To quantify absolute probabilistic accuracy under competing risks, the Brier Score serves as the equivalent metric, formulated directly against the Cumulative Incidence Function (CIF). Evaluated at a specific clinical time horizon $t$, the Brier Score calculates the mean squared difference between the model's predicted probability of the primary event $k$ occurring prior to time $t$, and the actual observed binary outcome.

Because survival data is subject to right censoring meaning the exact outcome is unknown for some patients at time $t$ the naive Brier Score is heavily biased. To resolve this, an Inverse Probability of Censoring Weighting (IPCW) term, $w_i(t)$, derived from the Kaplan-Meier censoring survival function $\hat{G}(\cdot)$, is dynamically applied. For a cohort of $N$ patients, the adjusted Brier Score is defined as:
$$\text{BS}(t) = \frac{1}{N} \sum_{i=1}^N w_i(t) \left( \mathbf{1}[T_i \le t, \Delta_i = k] - \hat{F}_k(t \mid \mathbf{x}_i) \right)^2$$
where $\mathbf{1}[\cdot]$ is the indicator function denoting whether the specific event $k$ was observed prior to time $t$. A lower Brier Score reflects superior calibration and tighter probabilistic agreement with empirical reality.

\subsection{Integrated Brier Score (IBS)}
While the point-in-time Brier Score is useful for specific clinical horizons (e.g., the median follow-up time), it is highly dependent on the chosen evaluation point and fails to capture the model's accuracy over the entire temporal trajectory. To resolve this, the Integrated Brier Score (IBS) is utilized as a global summary metric. 

The IBS computes the average of the Brier Score evaluated continuously over a prespecified temporal interval $[0, \tau]$:
$$\text{IBS} = \frac{1}{\tau} \int_0^{\tau} \text{BS}(t) \, dt$$
By integrating the predictive error across all valid time steps, the IBS collapses short-term and long-term probabilistic accuracy into a single scalar value. This metric is critical for objectively comparing disparate sampling strategies across the entire follow-up period.

\subsection{Time-Dependent Concordance Index (C-Index)}
While the Brier Score measures absolute probabilistic error, the Concordance Index (C-Index) quantifies the model's discriminative ability by evaluating how accurately the predicted relative risk scores rank the observed event times. In a clinical setting, discrimination answers the question: \textit{If Patient A experiences a fatal cardiac event before Patient B, did the model assign Patient A a higher risk score?}

Let $\eta_i$ and $\eta_j$ denote the predicted log-risk scores for a pair of subjects. A pair is considered concordant if the subject who fails earlier was assigned a higher risk score ($\eta_i > \eta_j$). Because right-censoring obscures the true temporal ordering for certain patients, an IPCW adjustment is integrated to ensure unbiased evaluation. To properly evaluate non-proportional hazards where risk predictions vary over time, we rely on the time-dependent concordance index proposed by Antolini et al. \cite{Antolini2005}. Adapted for time-dependent and competing risks scenarios, the estimator is defined as:
$$\widehat{C}_{\text{IPCW}} = \frac{ \sum_{i \neq j} \mathbf{1}(T_i < T_j, \Delta_i = k) \mathbf{1}(\eta_i > \eta_j) / \hat{G}^2(T_i) }{ \sum_{i \neq j} \mathbf{1}(T_i < T_j, \Delta_i = k) / \hat{G}^2(T_i) }$$
A value of $0.5$ indicates performance equivalent to random chance, whereas a value approaching $1.0$ proves strong clinical discriminative rank-ordering.

\subsection{Stratified Aalen-Johansen CIF Calibration}
Strong discrimination (a high C-Index) guarantees that the model correctly ranks patient risk, but it does not guarantee calibration. Calibration asks: \textit{If the model predicts a 20\% risk of an event for a group of patients, do approximately 20\% of those patients actually experience the event?} To prove that the absolute risk estimations output by the Fine-Gray loss are reliable, a Stratified CIF calibration diagnostic was developed.

Under this diagnostic, the validation cohort is stratified into discrete quantiles (e.g., Low, Medium, and High risk) based on the model's predicted relative risk scores. For each stratum, the mean predicted CIF curve is plotted against empirical reality. Because the standard Kaplan-Meier estimator is mathematically invalid in the presence of competing risks, the empirical reality is estimated strictly using the non-parametric Aalen-Johansen estimator \cite{aalen1978empirical}. If the model's predicted mean CIF curve closely shadows the step-wise Aalen-Johansen empirical curve across all risk strata, it proves that the neural network's absolute risk estimates are highly calibrated.

\subsection{CPU and Memory Profiling}
A fundamental hypothesis of this thesis is that sub-sampling the risk set via static, nested, or dynamic entropy-guided designs dramatically reduces the computational burden inherent to the full cohort Fine-Gray objective. Beyond predictive accuracy, this study explicitly evaluates the computational efficiency of these sampling regimes.

A continuous resource profiler is executed alongside the cross-validation training loop. Operating on an isolated parallel thread, the profiler polls the system to record absolute CPU utilization and Resident Set Size (RSS) memory consumption, ensuring that only the resources directly consumed by the model training process are captured. By measuring resource trajectories and total execution durations independently of external background noise, the study accurately quantifies the exact computational efficiency gains unlocked by the risk set approximation methods.

\subsection{Synthetic Data Verification (Pearson \& Spearman Correlation)}
A primary limitation of analyzing real-world medical data is the absence of an absolute "ground truth" for individual feature effects. While clinical literature can suggest trends, the exact mathematical coefficient for a specific covariate remains fundamentally unknown. Consequently, verifying whether an interpretable model has accurately isolated a specific feature's contribution is challenging.

To bypass this limitation, the synthetic cohort experiment incorporates Pearson and Spearman rank correlation coefficients. Because the synthetic dataset was generated with explicitly defined baseline hazards and exact true linear predictors, the ground truth is strictly known. The Pearson coefficient measures the degree of linear agreement between the model's predicted feature effects and the true parameters, while the Spearman coefficient evaluates the monotonic rank correlation. High correlation scores explicitly prove that the CRISPNAM architecture successfully recovers the true underlying biostatistical relationships without feature entanglement.


\chapter{Results and Discussion}
\section{Dataset Characteristics}

To robustly evaluate the proposed CRISPNAM architecture, including the Fine-Gray sub-distribution loss formulation and the integrated risk-set sampling strategies, this study utilized three distinct datasets: a large-scale synthetic dataset and two smaller real-world clinical benchmarks (PBC and Framingham).

The synthetic dataset ($N = 15,000$) was explicitly simulated to provide a rigorous mathematical environment with a high incidence of competing events. As illustrated in Figure \ref{fig:cif_synthetic}, the pronounced divergence between the standard Kaplan-Meier estimate and the true Cumulative Incidence Function (CIF) confirms that this dataset serves as an ideal stress test. It effectively demonstrates the architecture's capacity to correct the upward bias inherent in traditional single-event survival models. Furthermore, the large sample size ensures that the nested case-control and case-cohort sampling efficiencies are validated across a wide data range.

\begin{figure}[H]
    \centering
    \includegraphics[width=0.75\textwidth]{images/Data/fh.png} 
    \caption{Comparison of the Cumulative Incidence Function (CIF) and Kaplan-Meier (KM) estimate for the primary event in the synthetic dataset ($N = 15,000$). The substantial divergence illustrates the bias inherent in the KM estimator under severe competing risk conditions.}
    \label{fig:cif_synthetic}
\end{figure}

Subsequent evaluations were conducted on established, smaller-scale real-world benchmarks to assess the model's generalizability on standard, noisy clinical data. The Primary Biliary Cirrhosis (PBC) dataset, a foundational standard in survival analysis, comprises 1,945 observations. Within this cohort, 1,073 patients (55.2\%) were right-censored, 725 (37.3\%) experienced the primary event of interest, and 147 (7.6\%) experienced a competing event. While the lower proportion of competing events results in a less extreme deviation between the KM estimate and the CIF (Figure \ref{fig:cif_pbc}) compared to the synthetic data, the dataset remains a critical baseline for benchmarking clinical predictive performance against existing literature.

\begin{figure}[H]
    \centering
    \includegraphics[width=0.75\textwidth]{images/Data/pbc.png} 
    \caption{Cumulative Incidence Function (CIF) versus Kaplan-Meier (KM) estimate for the primary event in the real-world PBC dataset. The mild divergence reflects standard clinical event distributions.}
    \label{fig:cif_pbc}
\end{figure}

Finally, the semi-synthetic Framingham dataset was included as a secondary real-world evaluation point. Though significantly smaller in scale than the synthetic cohort and heavily censored (Figure \ref{fig:cif_framingham}), it contains persistent competing risk dynamics. Utilizing this combination of datasets ensures that the model's sampling mechanisms and loss functions are comprehensively validated across both mathematically extreme theoretical conditions and standard clinical realities.

\begin{figure}[H]
    \centering
    \includegraphics[width=0.75\textwidth]{images/Data/fh.png} 
    \caption{Cumulative Incidence Function (CIF) versus Kaplan-Meier (KM) estimate for the semi-synthetic Framingham dataset, demonstrating competing risk dynamics under high censoring conditions.}
    \label{fig:cif_framingham}
\end{figure}


\section{Optimization and Convergence}

Training deep survival architectures necessitates careful monitoring to prevent overfitting, particularly when optimizing complex sub-distribution hazard formulations across varying risk-set sampling designs. To ensure robust evaluation, all models were trained using a stratified 5-fold cross-validation strategy, ensuring consistent underlying event proportions across all splits. 

Given the extensive hyperparameter search space across the three datasets and four sampling strategies, plotting every trajectory is impractical. Instead, the loss curves presented reflect a single representative fold (Fold 1) utilizing the optimal hyperparameter configuration for each sampling type. The full cohort loss serves as the foundational baseline against which the sampling strategies are compared.

\subsection{Hyperparameter Stability and Selection}

The convergence stability of the CRISPNAM architecture was highly dependent on the hyperparameter configurations of the risk-set samplers. The optimal parameters were selected based on the highest cross-validated Concordance Index, revealing several consistent behavioral patterns across the datasets.

Nested Case-Control (NCC) Sampling we have found that the optimal number of controls was found to be $m = 5$. Increasing the number of controls did not yield significant predictive improvements but linearly increased computational overhead. This aligns with standard literature on deep survival sampling, notably Kvamme et al., who demonstrated that minimal control sets are sufficient to approximate the partial likelihood gradient effectively.

Static Case-Cohort (SCC) Sampling where the subcohort fraction of $\alpha = 0.5$ provided the most stable optimization trajectory. While lower fractions ($\alpha \in [0.2, 0.3]$) ultimately achieved similar terminal loss values, their convergence curves were highly jagged. Extreme subcohort compression ($\alpha < 0.2$) led to gradient collapse during early epochs. Thus, $\alpha = 0.5$ is recommended for maintaining gradient stability.

The Entropy Sampling makes use of a temperature scaling parameter ($\tau$)  that controls the greediness of the sampler. The optimal balance was found at $\tau = 0.5$. When the temperature was pushed toward maximum greediness ($\tau = 0.1$), the model severely overfit, as it exclusively sampled the highest-variance patients and failed to generalize to the broader distribution. Based on these empirical trajectories, a range of $\tau \in [0.5, 1.0]$ is recommended to adequately balance the exploitation of high-information risk sets with robust model generalizability.

\newpage

\subsection{Dataset Convergence Analysis}

As illustrated in Figure \ref{fig:loss_synthetic}, the Fine-Gray partial likelihood loss achieves textbook convergence on the large-scale synthetic dataset ($N=15,000$). The training and validation curves track almost identically across all four sampling configurations, verifying the structural and mathematical integrity of the gradients within the simulated competing risks environment. 

\begin{figure}[H]
    \centering
    \begin{subfigure}[b]{0.48\textwidth}
        \includegraphics[width=\textwidth]{images/Synth_results/FULL/LOSS/fold_1_loss_cause_1.png}
        \caption{Full Cohort}
    \end{subfigure}
    \hfill
    \begin{subfigure}[b]{0.48\textwidth}
        \includegraphics[width=\textwidth]{images/Synth_results/NCC/LOSS/fold_1_loss_cause_1_NCC_ncontrols_5.png}
        \caption{NCC ($m=5$)}
    \end{subfigure}
    \vskip\baselineskip
    \begin{subfigure}[b]{0.48\textwidth}
        \includegraphics[width=\textwidth]{images/Synth_results/SCC/LOSS/fold_1_loss_cause_1_SCC_subcohort_0.5.png}
        \caption{SCC ($\alpha=0.5$)}
    \end{subfigure}
    \hfill
    \begin{subfigure}[b]{0.48\textwidth}
        \includegraphics[width=\textwidth]{images/Synth_results/ENTROPY/LOSS/fold_1_loss_cause_1_Entropy_tau_0.5_compression_0.1.png}
        \caption{Entropy ($\tau=0.5$)}
    \end{subfigure}
    \caption{Loss trajectories for the synthetic dataset ($N=15,000$) under optimal hyperparameters, showing stable optimization in severe competing risk scenarios.}
    \label{fig:loss_synthetic}
\end{figure}

The optimization trajectories for the real-world PBC dataset (Figure \ref{fig:loss_pbc}) exhibit higher initial loss magnitudes, which is characteristic of standard, unscaled clinical data. However, the model successfully minimizes the sub-distribution hazard across all sampling configurations without catastrophic divergence. The validation loss consistently mirrors the training loss, proving that the network's internal regularization mechanisms prevent it from memorizing the sampled clinical risk-sets.

\begin{figure}[H]
    \centering
    \begin{subfigure}[b]{0.48\textwidth}
        \includegraphics[width=\textwidth]{images/pbc_results/FULL/LOSS/fold_1_loss_cause_1.png}
        \caption{Full Cohort}
    \end{subfigure}
    \hfill
    \begin{subfigure}[b]{0.48\textwidth}
        \includegraphics[width=\textwidth]{images/pbc_results/NCC/LOSS/fold_1_loss_cause_1_NCC_ncontrols_5.png}
        \caption{NCC ($m=5$)}
    \end{subfigure}
    \vskip\baselineskip
    \begin{subfigure}[b]{0.48\textwidth}
        \includegraphics[width=\textwidth]{images/pbc_results/SCC/LOSS/fold_1_loss_cause_1_SCC_subcohort_0.5.png}
        \caption{SCC ($\alpha=0.5$)}
    \end{subfigure}
    \hfill
    \begin{subfigure}[b]{0.48\textwidth}
        \includegraphics[width=\textwidth]{images/pbc_results/ENTROPY/LOSS/fold_1_loss_cause_1_Entropy_tau_0.5_compression_0.1.png}
        \caption{Entropy ($\tau=0.5$)}
    \end{subfigure}
    \caption{PBC dataset loss trajectories show stable convergence across all sampling methods despite early variance.}
    \label{fig:loss_pbc}
\end{figure}

Finally, the convergence trajectories for the semi-synthetic Framingham dataset (Figure \ref{fig:loss_fh}) present a unique visual scale. Because the initial unoptimized training loss is exceptionally high across all sampling configurations, the y-axis is heavily compressed. This visual compression makes the validation loss appear entirely flat at first glance. However, numerical inspection of the optimization logs confirms that the validation loss does indeed decrease steadily (e.g., dropping from approximately 6.6 to 6.1) as the network optimizes. 

This confirms that the model is actively learning and generalizing the sub-distribution hazard rather than stagnating. Importantly, across Full Cohort, NCC, SCC, and Entropy sampling, the training loss converges smoothly down to effectively meet the validation loss without crossing into severe overfitting territory. This further validates the robustness of the chosen hyperparameters when navigating heavily censored competing risks data.

\begin{figure}[H]
    \centering
    \begin{subfigure}[b]{0.48\textwidth}
        \includegraphics[width=\textwidth]{images/fh_results/FULL/LOSS/fold_1_loss_cause_1.png}
        \caption{Full Cohort}
    \end{subfigure}
    \hfill
    \begin{subfigure}[b]{0.48\textwidth}
        \includegraphics[width=\textwidth]{images/fh_results/NCC/LOSS/fold_1_loss_cause_1_NCC_ncontrols_5.png}
        \caption{NCC ($m=5$)}
    \end{subfigure}
    \vskip\baselineskip
    \begin{subfigure}[b]{0.48\textwidth}
        \includegraphics[width=\textwidth]{images/fh_results/SCC/LOSS/fold_1_loss_cause_1_SCC_subcohort_0.5.png}
        \caption{SCC ($\alpha=0.5$)}
    \end{subfigure}
    \hfill
    \begin{subfigure}[b]{0.48\textwidth}
        \includegraphics[width=\textwidth]{images/fh_results/ENTROPY/LOSS/fold_1_loss_cause_1_Entropy_tau_0.5_compression_0.1.png}
        \caption{Entropy ($\tau=0.5$)}
    \end{subfigure} 
    \caption{Framingham dataset loss trajectories demonstrate stable convergence and no overfitting, despite visual y-axis compression from high initial training loss.}
    \label{fig:loss_fh}
\end{figure}

\subsection{Computational Resource Profiling}

A primary motivation for implementing risk-set sampling strategies in deep survival architectures is the reduction of computational overhead. Calculating the exact sub-distribution partial likelihood for the full cohort becomes computationally intractable as dataset sizes scale. To quantify this efficiency gain, hardware resource utilization (CPU and Memory) and execution times were profiled during model training. 

Given that the computational patterns were consistent across all experiments, Figure \ref{fig:resource_profiling} visualizes the resource profiles exclusively for the large-scale synthetic dataset ($N=15,000$), which serves as the upper-bound stress test for the architecture. The smaller PBC and Framingham datasets exhibited identical behavioral patterns at lower absolute magnitudes.

\subsubsection{Resource Utilization and Execution Efficiency}

Across all sampling configurations, CPU utilization remained highly stable at approximately 50\% during the active training phase, indicating consistent dataloader throughput without severe processing bottlenecks. 

Memory utilization exhibited a characteristic pattern across all methods: an initial allocation spike during data ingestion and model initialization, followed by a stable plateau during batched gradient updates. 

    \paragraph{Full Cohort} The baseline full cohort calculation required approximately 725 MB during initialization and plateaued near 600 MB. However, it required the longest execution time, taking roughly 60 seconds per fold iteration.
    \paragraph{Static Case-Cohort (SCC) \& Nested Case-Control (NCC)} Both SCC ($\alpha=0.5$) and NCC ($m=5$) required slightly higher initial memory spikes (approx. 790 MB and 780 MB, respectively) due to the overhead of the sampling algorithms allocating the risk-sets in memory. However, by compressing the active risk pool, both methods reduced the execution time to approximately 45 seconds a 25\% reduction in compute time compared to the full cohort.
    \paragraph{Entropy Sampling} The entropy sampler ($\tau=0.5$) exhibited the highest initialization memory peak (approx. 830 MB), reflecting the upfront computational cost of calculating the entropy distribution across the risk-sets. Strikingly, once initialized, this strategy proved to be the most efficient, completing the execution in roughly 28 seconds. This represents a greater than 50\% reduction in training time compared to the baseline, proving that intelligently exploiting high-information risk sets drastically accelerates optimization.


Ultimately, these profiles confirm that while advanced sampling strategies require marginally more initialization memory to construct the active cohorts, they heavily optimize the overall computational pipeline by drastically reducing the time required for gradient convergence.

\section{Predictive performance and training time}
\subsection{Synthetic Dataset ($N=15,000$)}

The synthetic dataset ($N=15,000$) provides a robust mathematical stress test featuring a high incidence of competing events. This environment enables a clear evaluation of the CRISPNAM architecture's capacity to resolve competing risk bias while establishing the baseline predictive performance and computational overhead required for full partial likelihood calculations.



\begin{figure}[H]
    \centering
    \includegraphics[width=0.85\textwidth]{images/Synth_results/Pooled_C-Index.png}
    \caption{Pooled C-Index}
    \label{fig:synth_c_index}
\end{figure}

\subsubsection{Sampler: Full Cohort}
The full cohort model serves as the theoretical and empirical baseline. Because it utilizes the entire risk-set at every event time, it guarantees the most precise gradient updates at the expense of maximum computational complexity. As illustrated in Figure \ref{fig:synth_c_index} \ref{fig:synth_brier_score}, the full cohort achieved a pooled C-index of $0.850$ and an Brier Score of $0.107$. However, this mathematical precision comes with a significant computational penalty. The full cohort training required an average of $58.16 \pm 21.87$ seconds per fold, establishing the baseline latency that the risk-set sampling strategies aim to optimize.

\begin{figure}[H]
    \centering
    \includegraphics[width=0.85\textwidth]{images/Synth_results/Training_Duration_mean_std_error.png}
    \caption{Training Duration (Sec)}
    \label{fig:synth_training_duration}
\end{figure}

\subsubsection{Sampler: Nested Case-Control (NCC)}

The performance and computational cost of the NCC sampler are highly dependent on the number of controls ($m$). When restricted to $m=5$, the model experiences a slight degradation in discriminative capability (C-index: $0.835$) and calibration (Brier Score: $0.139$). Expanding the control pool to $m=100$ restores performance to baseline parity (C-index: $0.849$, Brier Score: $0.107$). Crucially, however, as the number of controls increases, both the training duration and the cross-fold variance scale linearly, with $m=100$ actually exceeding the training time of the full cohort (approaching 85 seconds on average). This behavior mirrors findings documented by Kvamme \cite{kvamme2019}. for deep single event Cox models, where arbitrarily large control sets introduce severe computational bottlenecks without proportional predictive gains. Furthermore, the massive standard deviation bars at higher control limits indicate that standard NCC struggles with execution stability across randomized data splits.

\begin{figure}[H]
    \centering
    \includegraphics[width=0.85\textwidth]{images/Synth_results/Pooled_Brier.png}
    \caption{Brier Score}
    \label{fig:synth_brier_score}
\end{figure}

\subsubsection{Sampler: Static Case-Cohort (SCC)}
The Static Case-Cohort sampler demonstrated high stability across varying subcohort fractions ($\alpha$). Even under extreme compression ($\alpha=0.05$), the model maintained a C-index of $0.842$ and an Brier Score of $0.110$. At the optimal configuration of $\alpha=0.5$, the metrics aligned near-perfectly with the full cohort (C-index: $0.849$, Brier Score: $0.107$) while compressing the training time to roughly 38 seconds. The tighter variance bars observed at $\alpha=0.5$ indicate that maintaining a static background cohort effectively stabilizes the model's convergence timing compared to the highly variable NCC approach.

\subsubsection{Sampler: Entropy}
The Entropy sampler provided the most significant finding of the synthetic evaluation. At optimal temperature scaling ($\tau \in [0.5, 1.0]$), the entropy sampler not only matched the discriminative performance of the Full Cohort baseline (C-index: $0.851$ at $\tau=0.5$) and maintained parity in calibration (Brier Score: $0.107$), but it achieved this in a fraction of the time. 

The training duration for the entropy sampler at $\tau=0.5$ was drastically reduced to approximately 29 seconds a 50\% reduction compared to the full cohort baseline. Furthermore, the cross-fold variance for training duration in the $\tau \in [0.5, 1.0]$ range is exceptionally small. This negligible standard deviation proves that the entropy formulation is highly deterministic; by intelligently weighting high information observations and filtering out redundant background noise.

\subsection{Real-World Clinical Benchmark: PBC Dataset ($N=1,945$)}

While the synthetic dataset validated the mathematical limits of the architecture under severe competing risk conditions, the Primary Biliary Cirrhosis (PBC) dataset serves to evaluate the model's robustness against standard clinical noise, unscaled covariates, and a smaller overall sample size ($N=1,945$). 

Because the dataset is significantly smaller, the baseline execution time is inherently lower. However, the performance metrics presented in Figure \ref{fig:pbc_brier_score}, \ref{fig:pbc_training_duration}, and \ref{fig:pbc_c_index} reveal that the sampling strategies react far more violently to hyperparameter tuning on noisy clinical data than they do on clean synthetic data.

\begin{figure}[H]
    \centering
    \includegraphics[width=0.85\textwidth]{images/pbc_results/Pooled_C-Index_PBC.png}
    \caption{Pooled C-Index}
    \label{fig:pbc_c_index}
\end{figure}

\subsubsection{Sampler: Full Cohort}
The full cohort model maintained strong discriminative capabilities on the clinical data, achieving a pooled C-index of 0.774 and an Integrated Brier Score (Brier Score) of 0.169. Due to the smaller dataset size, the computational burden was relatively light, averaging roughly 38 seconds per fold. Consequently, the primary objective of risk-set sampling on the PBC dataset shifts slightly: rather than strictly minimizing execution time, the goal is to prove that sampling does not catastrophically degrade the predictive integrity of a smaller cohort.


\begin{figure}[H]
    \centering
    \includegraphics[width=0.85\textwidth]{images/pbc_results/Pooled_Brier_PBC.png}
    \caption{Brier Score}
    \label{fig:pbc_brier_score}
\end{figure}

\subsubsection{Sampler: Nested Case-Control (NCC)}
The NCC sampler exhibited consistent behavior aligned with the synthetic findings. As the control parameter increased from $m=5$ to $m=100$, the predictive performance steadily climbed from a C-index of 0.740 to 0.773, reaching near-parity with the full cohort. However, the execution time scaled linearly alongside the control size. At $m=100$, NCC required approximately 42 seconds exceeding the duration of the full cohort model. This reinforces the finding that arbitrarily expanding control sets is computationally inefficient; the optimal balance for NCC on this dataset remains in the $m \in [20, 50]$ range.

\begin{figure}[H]
    \centering
    \includegraphics[width=0.85\textwidth]{images/pbc_results/Training_Duration_mean_std_error_PBC.png}
    \caption{Training Duration (Sec)}
    \label{fig:pbc_training_duration}
\end{figure}

\subsubsection{Sampler: Static Case-Cohort (SCC)}
The limitations of the Static Case-Cohort sampler were starkly exposed on the clinical data. At extreme subcohort compressions ($\alpha=0.05$), the model's discriminative ability collapsed entirely, plunging to a C-index of 0.571 only marginally better than random chance. The calibration similarly degraded, yielding an Brier Score of 0.276. This catastrophic failure indicates that when the overall dataset is small and noisy, heavily compressing the static background cohort removes critical variance necessary for stable gradient descent. The SCC model only approaches viability at $\alpha=0.5$ (C-index: 0.687), and even then, it fails to match the baseline.

\subsubsection{Sampler: Entropy}
The introduction of clinical noise forced a critical shift in the optimal temperature parameter ($\tau$) for the Entropy sampler. On the clean synthetic data, a highly greedy temperature ($\tau=0.5$) yielded optimal results. However, applying that same greedy parameter to the PBC data caused a significant performance drop (C-index: 0.665). 

This occurs because a low $\tau$ forces the sampler to exclusively focus on the highest-variance observations. In a noisy clinical setting, these "high-variance" patients are often extreme outliers. By over sampling outliers, the model fails to learn the general underlying survival distribution. 

The architecture recovers brilliantly when the temperature is relaxed to $\tau=1.0$ or $\tau=5.0$. By decreasing the greediness and introducing a flatter probability distribution over the risk-sets, the Entropy sampler recovers a C-index of 0.768 and an Brier Score of 0.171, reaching statistical parity with the full cohort. Furthermore, at this optimal temperature range, the execution time remains highly efficient (approximately 18-20 seconds), confirming that Entropy sampling remains the superior optimization strategy provided the temperature is scaled appropriately to the dataset's intrinsic noise profile.

\subsection{Real-World Semi-Synthetic: Framingham Dataset}

The Framingham dataset introduces a final, unique environmental challenge: it is heavily right-censored and features a lower overall event density. This scenario tests how the CRISPNAM architecture and its sampling mechanisms handle high uncertainty within the active risk pool. The  performance and computational profiling for this dataset are presented in Figure \ref{fig:fh_c_index}, \ref{fig:fh_brier_score}, and \ref{fig:fh_training_duration}.

\begin{figure}[H]
    \centering
    \includegraphics[width=0.85\textwidth]{images/fh_results/Pooled_C-Index_FH.png}
    \caption{Pooled C-Index}
    \label{fig:fh_c_index}
\end{figure}

\subsubsection{Sampler: Full Cohort}
Due to the heavy censoring and inherent noise, the Full Cohort baseline achieves a modest pooled C-index of 0.645 and a high Integrated Brier Score (Brier Score) of 0.670. Because the dataset is small, the entire risk matrix easily fits into GPU memory, allowing highly optimized, vectorized tensor operations. Consequently, the Full Cohort training is executed rapidly, averaging roughly 21 seconds per fold. 

\begin{figure}[H]
    \centering
    \includegraphics[width=0.85\textwidth]{images/fh_results/Pooled_Brier_FH.png}
    \caption{Integrated Brier Score}
    \label{fig:fh_brier_score}
\end{figure}


\subsubsection{Sampler: Nested Case-Control (NCC)}
The Framingham evaluation reveals a counter intuitive but critical behavior of the NCC sampler: it improves predictive performance but severely degrades computational efficiency. Across all control sizes ($m \in [5, 100]$), NCC consistently outperforms the Full Cohort baseline, achieving a C-index of up to 0.677 ($m=50$) and lowering the Brier Score to 0.644 ($m=5$). In heavily censored datasets, utilizing the full cohort forces the model to optimize against a massive pool of noisy, uncertain observations. NCC mitigates this by randomly discarding background noise, effectively acting as a data level regularizer that prevents overfitting.

However, the computational cost is inverted. At $m=5$, NCC requires approximately 34 seconds to train a 60\% increase over the Full Cohort. This is belived to occur, because the algorithmic overhead of dynamically sampling controls via CPU-bound loops outpaces the GPU compute time saved by the reduced matrix size. This confirms that on small datasets, the vectorization benefits of the Full Cohort outweigh the data-compression benefits of NCC.

\begin{figure}[H]
    \centering
    \includegraphics[width=0.85\textwidth]{images/fh_results/Training_Duration_mean_std_error_FH.png}
    \caption{Training Duration (Sec)}
    \label{fig:fh_training_duration}
\end{figure}

\subsubsection{Sampler: Static Case-Cohort (SCC)}
Similar to the PBC dataset, SCC struggles with stability when the background cohort is too small. At $\alpha=0.05$, the C-index collapses to 0.588. While it stabilizes at $\alpha=0.5$ (C-index: 0.637), it fails to capture the regularization benefits seen in NCC and suffers from high calibration error (Brier Score: 0.681). Its execution time remains comparable to the baseline, offering no distinct advantage in this highly censored environment.

\subsubsection{Sampler: Entropy}
The Entropy sampler reinforces the tuning dynamics observed in the PBC dataset. At an aggressively greedy temperature ($\tau=0.1$), the model overfits to the highest variance do to the  clinical outliers, resulting in a collapsed C-index of 0.575 and a severe Brier Score spike (0.863). 

However, when properly tuned to $\tau=5.0$, the entropy mechanism successfully filters the heavy censoring noise. It achieves a C-index of 0.663 and an Brier Score of 0.655, outperforming the Full Cohort baseline. Furthermore, unlike NCC, the algorithmic implementation of the Entropy sampler is highly efficient at optimal temperatures; at $\tau=0.5$, it executes in roughly 17 seconds, successfully improving upon both the predictive and computational benchmarks of the Full Cohort.


\section{Model Calibration and Discrimination}

While the previous section established the predictive performance using the pooled Concordance Index and Brier Score, these metrics primarily evaluate general rank-ordering and time-averaged error. To fully validate the CRISPNAM architecture, it is necessary to examine its calibration its ability to predict absolute event probabilities accurately over time and its capacity to recover true underlying biostatistical relationships.

\subsection{Ground Truth Risk Recovery}

A fundamental limitation of real-world clinical datasets is that the true feature effects and baseline hazards are strictly unknown. The synthetic dataset ($N=15,000$) bypasses this limitation, allowing for a direct comparison between the model's predicted relative risk score ($\hat{\eta}$) and the mathematically defined True Linear Predictor.

Figure \ref{fig:synth_truth_recovery} illustrates this relationship for the Full Cohort baseline. The architecture successfully isolates the true hazard contributions, achieving a Spearman rank correlation of 0.833 and a Pearson linear correlation of 0.843 for the primary event of interest. This strong linear and monotonic agreement proves that the deep neural additive framework accurately recovers the true biostatistical feature effects without entanglement, even under severe competing risk conditions.

\begin{figure}[htbp]
    \centering
    \includegraphics[width=0.85\textwidth]{images/Synth_results/FULL/Metrics/model_diagnostics_cause_1.png}
    \caption{Model diagnostics for the Full Cohort on the synthetic dataset. The left panel demonstrates strong ground truth risk recovery (Pearson $r=0.843$). The right panel confirms absolute probabilistic accuracy, with the model's predicted Cumulative Incidence Functions closely tracking the empirical Aalen-Johansen estimates across three distinct risk strata.}
    \label{fig:synth_truth_recovery}
\end{figure}

\subsection{Cumulative Incidence Calibration}

Beyond rank-ordering patients correctly, a survival model must be highly calibrated; if the network predicts a 40\% cumulative incidence for a specific subgroup at a specific time horizon, empirical reality must reflect that exact proportion. 

To validate this absolute accuracy, the synthetic cohort was partitioned into three distinct risk strata (Low, Medium, and High) based on the network's predicted scores. As shown in the right panel of Figure \ref{fig:synth_truth_recovery}, the predicted mean CIF curves for all three strata effectively overlay the non-parametric Empirical Aalen-Johansen estimators. There is no visible divergence, probability collapse, or inflation at the temporal bounds, confirming that the unmodified Fine-Gray objective function is deeply calibrated.

\begin{figure}[htbp]
    \centering
    \includegraphics[width=0.85\textwidth]{images/Synth_results/FULL/Metrics/subject_curves_cause_1_cif.png}
    \caption{Individual patient Cumulative Incidence trajectories compared against the population Aalen-Johansen baseline. The left panel illustrates extreme risk separation by plotting the four highest and four lowest-risk subjects. The right panel displays the individual trajectories of five randomly sampled subjects across the risk spectrum. The model successfully generates highly individualized, monotonic probability curves that accurately reflect subject-specific risk dynamics without violating the subdistribution boundaries.}
    \label{fig:synth_individual_cif}
\end{figure}

\subsection{Entropy Sampler Calibration and Risk Recovery}

A core hypothesis of this study is that intelligent risk-set sampling can drastically reduce computational overhead without corrupting the model's calibration or its ability to recover true biostatistical relationships. To evaluate this, the dynamic Entropy sampling strategy was benchmarked against the Full Cohort baseline.

Remarkably, the diagnostic evaluations reveal that the Entropy sampler does not merely match the baseline; at optimal temperature scaling ($\tau=0.5$), it strictly outperforms it in ground truth recovery. As illustrated in Figure \ref{fig:entropy_truth_recovery}, the Entropy sampler achieves a Pearson linear correlation of 0.866 and a Spearman rank correlation of 0.857. This is a measurable improvement over the Full Cohort's respective scores of 0.843 and 0.833. By dynamically prioritizing high-variance, high-information subjects and aggressively down-sampling dense, redundant patient clusters, the entropy mechanism effectively acts as a data-level regularizer, allowing the neural additive framework to isolate a cleaner underlying hazard signal.

Furthermore, as the temperature parameter $\tau$ is increased forcing the probability distribution to become flatter and less greedy the model's recovery metrics gracefully degrade back toward the Full Cohort baseline. At $\tau=1.0$, the Pearson correlation drops to 0.853, and at $\tau=5.0$, it falls to 0.839. This confirms the expected theoretical behavior: as the dynamic sampler approaches a uniform random sample, its regularizing benefits diminish.

\begin{figure}[htbp]
    \centering
    \includegraphics[width=0.85\textwidth]{images/Synth_results/ENTROPY/Metrics/model_diagnostics_cause_1_Entropy_tau_0.5_compression_0.1.png}
    \caption{Model diagnostics for the optimal Entropy Sampler ($\tau=0.5$). The left panel shows superior risk score recovery (Pearson $r=0.866$) compared to the full risk-set. The right panel proves that Inverse Probability Weighting (IPW) successfully corrects selection bias, yielding highly calibrated CIF curves across all risk strata.}
    \label{fig:entropy_truth_recovery}
\end{figure}

Crucially, this aggressive sub-sampling does not destroy the absolute probability estimates. A known vulnerability of risk-set sampling is the introduction of severe selection bias. However, the right panel of Figure \ref{fig:entropy_truth_recovery} demonstrates that the Inverse Probability Weighting (IPW) integration perfectly neutralizes this bias. The predicted Stratified CIF curves cleanly track the empirical Aalen-Johansen estimates without probability inflation.

\begin{figure}[htbp]
    \centering
    \includegraphics[width=0.85\textwidth]{images/Synth_results/ENTROPY/Metrics/subject_curves_cause_1_cif_Entropy_tau_0.5_compression_0.1.png}
    \caption{Individual patient Cumulative Incidence trajectories generated via the Entropy sampler ($\tau=0.5$). The left panel isolates the four highest and lowest risk subjects, while the right panel displays five randomly sampled individuals. The dynamic sampling algorithm successfully preserves individual monotonicity and absolute scale without violating subdistribution limits.}
    \label{fig:entropy_individual_cif}
\end{figure}

This calibration is further maintained at the micro-level. Figure \ref{fig:entropy_individual_cif} displays the individual subject CIF trajectories generated by the $\tau=0.5$ Entropy model. Despite being trained on a heavily fractionated approximation of the cohort, the network continues to output stable, monotonic, and bounded incidence curves that appropriately scale to specific patient covariate profiles.


\subsection{Nested Case-Control (NCC) Calibration and Risk Recovery}

Following the evaluation of the dynamic Entropy sampler, the Nested Case-Control (NCC) architecture was assessed to determine how rigid, uniform control sampling impacts absolute calibration and feature recovery. As established in the macro-performance analysis, NCC is highly sensitive to the defined control quota ($m$).

When provided with a large control pool ($m=100$), the NCC sampler successfully mimics the full cohort. As shown in Figure \ref{fig:ncc_truth_recovery_100}, this configuration achieves a Pearson correlation of 0.856 and a Spearman rank correlation of 0.847, marginally outperforming the Full Cohort baseline by filtering out a portion of background noise. The Inverse Probability Weighting (IPW) mechanism functions effectively at this scale, ensuring the predicted stratified CIF curves accurately align with the empirical Aalen-Johansen estimates without significant bias.

\begin{figure}[htbp]
    \centering
    \includegraphics[width=0.85\textwidth]{images/Synth_results/NCC/Metrics/model_diagnostics_cause_1_NCC_ncontrols_100.png}
    \caption{Model diagnostics for NCC with an optimal control pool ($m=100$). The architecture successfully recovers the ground truth linear predictor (Pearson $r=0.856$) and produces highly calibrated CIF curves across all risk strata, closely matching the empirical baseline.}
    \label{fig:ncc_truth_recovery_100}
\end{figure}

However, evaluating the boundary conditions of the NCC sampler reveals a critical vulnerability when the active risk-set is overly compressed. Figure \ref{fig:ncc_truth_recovery_5} illustrates the catastrophic breakdown of the model when restricted to just $m=5$ controls per case. While the linear correlation degrades to 0.806, the most severe failure occurs in the absolute probability calibration. 

\begin{figure}[htbp]
    \centering
    \includegraphics[width=0.85\textwidth]{images/Synth_results/NCC/Metrics/model_diagnostics_cause_1_NCC_ncontrols_5.png}
    \caption{Model diagnostics for a heavily constrained NCC sampler ($m=5$). While rank-ordering partially survives, absolute probability calibration collapses completely. The high-risk stratum (Stratum 3, green) terminates abruptly, and the model severely underestimates the true cumulative incidence, visually explaining the Brier Score degradation observed in earlier evaluations.}
    \label{fig:ncc_truth_recovery_5}
\end{figure}

The Stratified CIF calibration plot for $m=5$ demonstrates a total structural collapse. The high-risk curve (Stratum 3) terminates prematurely near $t=100$, and all strata drastically underestimate the true Aalen-Johansen incidence. This occurs because starving the neural network of valid controls artificially inflates the perceived baseline hazard during mini-batch optimization, distorting the IPW scaling. This visual diagnostic directly explains why the Brier Score spiked to 0.139 in the previous section: without sufficient controls, standard NCC cannot maintain absolute probabilistic accuracy over time.

\begin{figure}[htbp]
    \centering
    \includegraphics[width=0.85\textwidth]{images/Synth_results/NCC/Metrics/subject_curves_cause_1_cif_NCC_ncontrols_100.png}
    \caption{Individual patient Cumulative Incidence trajectories generated via NCC ($m=100$). When provided with sufficient controls, the model retains the ability to generate stable, individualized subdistribution probabilities that scale accurately across the temporal domain.}
    \label{fig:ncc_individual_cif}
\end{figure}

Conversely, when operating at the optimal $m=100$ configuration, the micro-level calibration remains intact. Figure \ref{fig:ncc_individual_cif} confirms that the individual patient trajectories are stable, monotonic, and effectively separate high-risk and low-risk subjects while respecting the mathematical boundaries of the subdistribution hazard framework.

\subsection{Static Case-Cohort (SCC) Calibration and Risk Recovery}

The Static Case-Cohort (SCC) architecture offers a fundamentally different sampling philosophy, relying on a fixed background subcohort rather than time-matched controls. While this approach provides stable optimization across iterations, its ability to maintain absolute calibration is highly dependent on the retention fraction ($\alpha$).

When provided with a robust background cohort ($\alpha=0.5$), the SCC sampler demonstrates strong structural integrity. As shown in Figure \ref{fig:scc_truth_recovery_05}, this configuration recovers the ground truth linear predictor with a Pearson correlation of 0.849 and a Spearman correlation of 0.840, slightly exceeding the Full Cohort baseline. Importantly, the IPW integration successfully scales the fixed background cohort to match the population distribution, resulting in Stratified CIF curves that accurately align with the empirical Aalen-Johansen estimates across all three risk strata.

\begin{figure}[htbp]
    \centering
    \includegraphics[width=0.85\textwidth]{images/Synth_results/SCC/Metrics/model_diagnostics_cause_1_SCC_subcohort_0.5.png}
    \caption{Model diagnostics for SCC with an optimal subcohort fraction ($\alpha=0.5$). The model successfully recovers the ground truth hazard (Pearson $r=0.849$) and maintains highly calibrated absolute probability estimates across all risk strata.}
    \label{fig:scc_truth_recovery_05}
\end{figure}

However, similar to the NCC architecture, aggressively compressing the risk-set exposes severe methodological limits. When the background subcohort is reduced to just 5\% of the total population ($\alpha=0.05$), the model's calibration begins to fracture, as illustrated in Figure \ref{fig:scc_truth_recovery_05_collapse}. 

\begin{figure}[htbp]
    \centering
    \includegraphics[width=0.85\textwidth]{images/Synth_results/SCC/Metrics/model_diagnostics_cause_1_SCC_subcohort_0.05.png}
    \caption{Model diagnostics for a heavily compressed SCC sampler ($\alpha=0.05$). While rank-ordering correlation remains relatively high (Pearson $r=0.835$), the absolute calibration degrades. The Stratified CIF curves begin to diverge from the empirical baseline, particularly in the lower risk strata, indicating a breakdown in IPW scaling due to cohort sparsity.}
    \label{fig:scc_truth_recovery_05_collapse}
\end{figure}

Interestingly, the nature of the SCC collapse differs from the NCC collapse. While the severely constrained NCC model ($m=5$) completely lost structural integrity and terminated early, the heavily constrained SCC model ($\alpha=0.05$) maintains its basic shape and still achieves a reasonable Pearson correlation of 0.835. The failure occurs entirely in the absolute calibration: because the static subcohort is too sparse to accurately represent the population's covariate density, the IPW scaling becomes skewed. Consequently, the predicted Stratified CIF curves systematically pull away from the true empirical baseline, particularly for the low and medium risk groups.

\begin{figure}[htbp]
    \centering
    \includegraphics[width=0.85\textwidth]{images/Synth_results/SCC/Metrics/subject_curves_cause_1_cif_SCC_subcohort_0.5.png}
    \caption{Individual patient Cumulative Incidence trajectories generated via SCC ($\alpha=0.5$). The fixed background cohort provides sufficient gradient stability to generate smooth, individualized subdistribution probability curves that properly scale to patient risk.}
    \label{fig:scc_individual_cif}
\end{figure}

Provided the subcohort fraction remains sufficiently large (e.g., $\alpha=0.5$), the SCC architecture remains highly reliable at the individual patient level. Figure \ref{fig:scc_individual_cif} demonstrates that the optimal SCC configuration successfully generates individualized, monotonic CIF trajectories that respect the subdistribution boundaries and cleanly separate high and low-risk clinical profiles.


\section{Real-World Model Calibration and Discrimination}

While the synthetic dataset provided a mathematically pure environment to test absolute boundaries and feature recovery, standard clinical data introduces high variance, unscaled covariates, and inherent observational noise. To validate the CRISPNAM architecture's generalizability, the calibration of the risk-set sampling mechanisms was evaluated against the real-world Primary Biliary Cirrhosis (PBC) benchmark.

Because the ground truth linear predictor is fundamentally unknown in real-world clinical data, the evaluation relies exclusively on empirical calibration diagnostics. Specifically, the model's predicted Cumulative Incidence Functions (CIF) are plotted against the non-parametric Aalen-Johansen estimators, evaluated both at a stratified group level and via marginal calibration gaps.

\subsection{Real-World Clinical Benchmark: PBC Dataset}

\subsubsection{Full Cohort Baseline}

The Full Cohort configuration serves as the empirical baseline for the PBC dataset. By computing the partial likelihood over the entire available risk set at every event time, it establishes the maximum achievable calibration given the intrinsic noise of the data.

Figure \ref{fig:pbc_full_calibration} illustrates this baseline calibration. The left panel displays the Stratified CIF Calibration; the clinical cohort was partitioned into three distinct risk strata (Low, Medium, and High) based on the network's predicted relative risk scores. The model's predicted mean incidence (solid lines) tracks the empirical Aalen-Johansen estimates (dashed lines) closely across all three risk levels. While there is minor localized deviation a standard artifact of noisy, real-world medical data the overall probabilistic scale is accurate, and there is no severe divergence or probability collapse over the 5,000-day observational window.

\begin{figure}[htbp]
    \centering
    \begin{subfigure}[b]{0.48\textwidth}
        \includegraphics[width=\textwidth]{images/pbc_results/FULL/Metrics/cif_calibration_cause_1.png}
        \caption{Stratified CIF Calibration}
    \end{subfigure}
    \hfill
    \begin{subfigure}[b]{0.48\textwidth}
        \includegraphics[width=\textwidth]{images/pbc_results/FULL/Metrics/subject_curves_cause_1_cif.png}
        \caption{Marginal Calibration and Discrimination}
    \end{subfigure}
    \caption{Calibration diagnostics for the Full Cohort baseline on the PBC dataset. The Stratified CIF curves (left) track empirical reality across risk bounds. The Marginal Calibration plot (right) demonstrates a tight agreement between the mean model prediction and the population baseline, confirming robust absolute risk estimation without selection bias.}
    \label{fig:pbc_full_calibration}
\end{figure}

The right panel of Figure \ref{fig:pbc_full_calibration} introduces the Marginal Calibration diagnostic. This metric compares the mean of all individual patient predictions (orange line) directly against the overall population Aalen-Johansen baseline (dashed black line). The shaded "Calibration Gap" remains exceptionally narrow throughout the entire timeline, confirming that the Full Cohort model yields unbiased absolute risk estimates without systematically over  or under-predicting the incidence of the primary event.

\subsubsection{Nested Case-Control (NCC)}

The calibration of the Nested Case-Control (NCC) architecture on the PBC dataset heavily mirrors the failure states observed in the synthetic environment, proving that the algorithm's vulnerability to extreme risk-set compression is a structural limitation, not an artifact of data simulation.

When provided with a robust control pool ($m=100$), the NCC sampler successfully recovers the baseline calibration. As seen in Figure \ref{fig:pbc_ncc_100_calibration}, both the Stratified CIF and Marginal Calibration plots look nearly identical to the Full Cohort. The IPW mechanism functions as intended, neutralizing the sampling bias and preventing the Marginal Calibration Gap from expanding.

\begin{figure}[htbp]
    \centering
    \begin{subfigure}[b]{0.48\textwidth}
        \includegraphics[width=\textwidth]{images/pbc_results/NCC/Metrics/cif_calibration_cause_1_NCC_ncontrols_100.png}
        \caption{Stratified CIF Calibration ($m=100$)}
    \end{subfigure}
    \hfill
    \begin{subfigure}[b]{0.48\textwidth}
        \includegraphics[width=\textwidth]{images/pbc_results/NCC/Metrics/subject_curves_cause_1_cif_NCC_ncontrols_100.png}
        \caption{Marginal Calibration ($m=100$)}
    \end{subfigure}
    \caption{Calibration diagnostics for NCC with a large control pool ($m=100$). The model successfully replicates the Full Cohort baseline, maintaining tight Stratified CIF tracking and a minimal Marginal Calibration Gap.}
    \label{fig:pbc_ncc_100_calibration}
\end{figure}

However, when the NCC sampler is aggressively constrained to just $m=5$ controls per case, the absolute calibration degrades significantly, despite the dataset being much smaller than the synthetic cohort. Figure \ref{fig:pbc_ncc_5_calibration} illustrates this breakdown. In the left panel, the high-risk stratum (green line) systematically underestimates the empirical reality, pulling away from the Aalen-Johansen estimator early in the timeline. 

\begin{figure}[htbp]
    \centering
    \begin{subfigure}[b]{0.48\textwidth}
        \includegraphics[width=\textwidth]{images/pbc_results/NCC/Metrics/cif_calibration_cause_1_NCC_ncontrols_5.png}
        \caption{Stratified CIF Calibration ($m=5$)}
    \end{subfigure}
    \hfill
    \begin{subfigure}[b]{0.48\textwidth}
        \includegraphics[width=\textwidth]{images/pbc_results/NCC/Metrics/subject_curves_cause_1_cif_NCC_ncontrols_5.png}
        \caption{Marginal Calibration ($m=5$)}
    \end{subfigure}
    \caption{Calibration diagnostics for a heavily constrained NCC sampler ($m=5$). The extreme sparsity of the control pool disrupts the IPW scaling, causing the high-risk stratum to underestimate the true incidence. Consequently, the Marginal Calibration Gap (right panel) expands severely over time as the mean model prediction detaches from the population baseline.}
    \label{fig:pbc_ncc_5_calibration}
\end{figure}

The consequence of this failure is visually stark in the Marginal Calibration plot (right panel). The mean model prediction severely detaches from the population baseline, creating a massive, expanding Calibration Gap. By starving the neural network of valid controls at each event time, the IPW formulation loses its stabilizing context, causing the model to systematically underestimate the baseline hazard of the clinical population. This definitively proves that while ultra-low control NCC ($m=5$) provides maximum computational speed, it sacrifices the absolute probabilistic integrity required for reliable clinical deployment.
\subsubsection{Static Case-Cohort (SCC)}

The Static Case-Cohort (SCC) architecture relies on a fixed background subcohort, theoretically providing a more stable optimization target than the dynamically changing risk sets of the NCC approach. However, applying this method to the small, noisy Primary Biliary Cirrhosis (PBC) dataset reveals significant vulnerabilities to cohort sparsity.

Even under an optimal, low-compression configuration ($\alpha=0.5$), where 50\% of the population is retained as the background cohort, the SCC sampler struggles to match the baseline calibration. As shown in Figure \ref{fig:pbc_scc_50_calibration}, while the general shape of the Stratified CIF curves is preserved, there is noticeable deviation across all risk strata compared to the empirical Aalen-Johansen estimates.

\begin{figure}[htbp]
    \centering
    \begin{subfigure}[b]{0.48\textwidth}
        \includegraphics[width=\textwidth]{images/pbc_results/SCC/Metrics/cif_calibration_cause_1_SCC_subcohort_0.5.png}
        \caption{Stratified CIF Calibration ($\alpha=0.5$)}
    \end{subfigure}
    \hfill
    \begin{subfigure}[b]{0.48\textwidth}
        \includegraphics[width=\textwidth]{images/pbc_results/SCC/Metrics/subject_curves_cause_1_cif_SCC_subcohort_0.5.png}
        \caption{Marginal Calibration ($\alpha=0.5$)}
    \end{subfigure}
    \caption{Calibration diagnostics for SCC with an optimal subcohort fraction ($\alpha=0.5$). Even with 50\% of the data retained, the model exhibits localized deviations in the Stratified CIF curves and systematic underestimation in the Marginal Calibration Gap.}
    \label{fig:pbc_scc_50_calibration}
\end{figure}

The consequence of these localized deviations is apparent in the Marginal Calibration plot (right panel of Figure \ref{fig:pbc_scc_50_calibration}). The mean model prediction systematically underestimates the population baseline, creating a persistent and expanding calibration gap. Unlike the synthetic dataset, which contained 15,000 observations, the PBC dataset is highly constrained. Therefore, discarding even 50\% of the background cohort removes critical covariate density required for accurate Inverse Probability Weighting (IPW), resulting in a dampened absolute risk profile.

This vulnerability becomes catastrophic under heavy compression. Figure \ref{fig:pbc_scc_05_calibration} illustrates the model's behavior when the background subcohort is reduced to just 5\% ($\alpha=0.05$). 

\begin{figure}[htbp]
    \centering
    \begin{subfigure}[b]{0.48\textwidth}
        \includegraphics[width=\textwidth]{images/pbc_results/SCC/Metrics/cif_calibration_cause_1_SCC_subcohort_0.05.png}
        \caption{Stratified CIF Calibration ($\alpha=0.05$)}
    \end{subfigure}
    \hfill
    \begin{subfigure}[b]{0.48\textwidth}
        \includegraphics[width=\textwidth]{images/pbc_results/SCC/Metrics/subject_curves_cause_1_cif_SCC_subcohort_0.05.png}
        \caption{Marginal Calibration ($\alpha=0.05$)}
    \end{subfigure}
    \caption{Calibration diagnostics for a heavily compressed SCC sampler ($\alpha=0.05$). The extreme sparsity of the static subcohort entirely destroys the IPW scaling mechanism. The model fails to recognize the baseline hazard, generating effectively flat predictions and resulting in total calibration collapse.}
    \label{fig:pbc_scc_05_calibration}
\end{figure}

At $\alpha=0.05$, the IPW mechanism experiences a total structural collapse. The Stratified CIF curves (left panel) are effectively crushed to zero, failing entirely to track the empirical risk progression. Because the static subcohort is too sparse to represent the background hazard of the clinical population, the network optimizes against a distorted, near-zero baseline. This is confirmed by the Marginal Calibration plot (right panel), where the mean model prediction remains flat, creating a massive calibration gap that encompasses nearly the entire population incidence. 

This finding demonstrates that while static subcohorts offer computational stability in large datasets, they are mathematically fragile when applied to smaller clinical trials, as they risk destroying the absolute probability calibration of the subdistribution framework.

\subsubsection{Entropy Sampler}

The application of the dynamic Entropy sampler to the Primary Biliary Cirrhosis (PBC) dataset reveals a critical insight regarding hyperparameter tuning in the presence of real-world clinical noise. In the mathematically pure synthetic dataset, a highly greedy temperature scaling ($\tau=0.5$) acted as a powerful regularizer, filtering out redundant patients and improving risk recovery. However, applying this exact same greedy configuration to a small, noisy clinical dataset results in a severe calibration failure.

Figure \ref{fig:pbc_entropy_05_calibration} illustrates the collapse of the Entropy sampler at $\tau=0.5$. Because the PBC dataset is inherently noisy, forcing the sampler to hyper-fixate on the absolute highest-variance observations causes it to overfit to clinical outliers. This extreme sampling skew destroys the integrity of the Inverse Probability Weighting (IPW) mechanism. 

\begin{figure}[htbp]
    \centering
    \begin{subfigure}[b]{0.48\textwidth}
        \includegraphics[width=\textwidth]{images/pbc_results/Entropy/Metrics/cif_calibration_cause_1_Entropy_tau_0.5_compression_0.1.png}
        \caption{Stratified CIF Calibration ($\tau=0.5$)}
    \end{subfigure}
    \hfill
    \begin{subfigure}[b]{0.48\textwidth}
        \includegraphics[width=\textwidth]{images/pbc_results/Entropy/Metrics/subject_curves_cause_1_cif_Entropy_tau_0.5_compression_0.1.png}
        \caption{Marginal Calibration ($\tau=0.5$)}
    \end{subfigure}
    \caption{Calibration diagnostics for a highly greedy Entropy sampler ($\tau=0.5$). Hyper-fixating on high-variance observations in a noisy clinical dataset corrupts the IPW scaling. The model systematically underestimates absolute risk across all strata (left) and exhibits a massive Marginal Calibration Gap (right).}
    \label{fig:pbc_entropy_05_calibration}
\end{figure}

As seen in the left panel of Figure \ref{fig:pbc_entropy_05_calibration}, the predicted Stratified CIF curves systematically pull away and underestimate the true Aalen-Johansen incidence across all risk groups. Consequently, the Marginal Calibration plot (right panel) exhibits a massive, expanding calibration gap, as the mean model prediction entirely fails to capture the true baseline hazard of the population.

To correct this failure, the entropy of the sampling distribution must be artificially increased. By raising the temperature parameter to $\tau=5.0$, the sampling distribution becomes flatter and less greedy, forcing the model to retain a wider, more representative spread of the clinical cohort while still prioritizing information-dense observations. 

\begin{figure}[htbp]
    \centering
    \begin{subfigure}[b]{0.48\textwidth}
        \includegraphics[width=\textwidth]{images/pbc_results/Entropy/Metrics/cif_calibration_cause_1_Entropy_tau_5_compression_0.1.png}
        \caption{Stratified CIF Calibration ($\tau=5.0$)}
    \end{subfigure}
    \hfill
    \begin{subfigure}[b]{0.48\textwidth}
        \includegraphics[width=\textwidth]{images/pbc_results/Entropy/Metrics/subject_curves_cause_1_cif_Entropy_tau_5_compression_0.1.png}
        \caption{Marginal Calibration ($\tau=5.0$)}
    \end{subfigure}
    \caption{Calibration diagnostics for the optimal Entropy sampler ($\tau=5.0$). By softening the greedy sampling distribution, the IPW mechanism is restored. The model successfully recovers the Full Cohort baseline calibration, accurately tracking both the Stratified CIF curves and the population marginal average.}
    \label{fig:pbc_entropy_5_calibration}
\end{figure}

As demonstrated in Figure \ref{fig:pbc_entropy_5_calibration}, this hyperparameter shift completely restores the model's absolute calibration. The Stratified CIF curves (left panel) closely track the empirical non-parametric estimates, and the Marginal Calibration Gap (right panel) narrows significantly, matching the structural integrity of the Full Cohort baseline. 

This behavior highlights the adaptive strength of the Entropy formulation: the temperature parameter $\tau$ acts as a tunable dial between computational compression and regularization. For large, clean datasets, it can be tuned down to aggressively compress the risk-set; for small, noisy clinical cohorts, it can be tuned up to preserve absolute probabilistic calibration without resorting to the brute-force computational overhead of a Full Cohort partial likelihood.

\subsection{Real-World Semi-Synthetic: Framingham Dataset}

To conclude the calibration analysis, the sampling architectures were evaluated on the Framingham heart study dataset. Unlike the highly distinct risk profiles of the synthetic data or the dense event rate of the PBC cohort, the Framingham dataset introduces a different biostatistical challenge: severe right-censoring. With a primary event incidence of only 15.6\%, the majority of the cohort is censored, creating an inherently sparse risk landscape over the long-term observational horizon.

\subsubsection{Full Cohort Baseline}

The Full Cohort configuration once again serves as the gold-standard baseline. By leveraging the entire available risk-set, it maximizes the amount of information extracted from the heavily censored data. 

As illustrated in Figure \ref{fig:fh_full_calibration}, the Full Cohort model handles the high censoring rate exceptionally well. The Stratified CIF Calibration (left panel) demonstrates that the model's predicted incidence correctly scales with empirical reality, even for the highest risk stratum, without experiencing probability inflation at the tail end of the survival curve. Furthermore, the Marginal Calibration diagnostic (right panel) shows near-perfect alignment between the mean model prediction and the population Aalen-Johansen baseline, with no visible calibration gap over the 3,500-day timeline.

\begin{figure}[htbp]
    \centering
    \begin{subfigure}[b]{0.48\textwidth}
        \includegraphics[width=\textwidth]{images/fh_results/FULL/Metrics/cif_calibration_cause_1.png}
        \caption{Stratified CIF Calibration}
    \end{subfigure}
    \hfill
    \begin{subfigure}[b]{0.48\textwidth}
        \includegraphics[width=\textwidth]{images/fh_results/FULL/Metrics/subject_curves_cause_1_cif.png}
        \caption{Marginal Calibration and Discrimination}
    \end{subfigure}
    \caption{Calibration diagnostics for the Full Cohort baseline on the heavily censored Framingham dataset. The model maintains tight Stratified CIF tracking (left) and exhibits no marginal calibration gap (right), successfully isolating the baseline hazard despite low event density.}
    \label{fig:fh_full_calibration}
\end{figure}

\subsubsection{Nested Case-Control (NCC)}

The high censoring rate in the Framingham dataset exposes a unique interaction with the Nested Case-Control (NCC) sampler. In previous datasets, heavy risk-set compression ($m=5$) caused the IPW formulation to break down, resulting in a massive underestimation of absolute risk. However, as shown in Figure \ref{fig:fh_ncc_calibration}, the NCC architecture is remarkably resilient on this specific dataset, even at extreme sparsity.

\begin{figure}[htbp]
    \centering
    \begin{subfigure}[b]{0.48\textwidth}
        \includegraphics[width=\textwidth]{images/fh_results/NCC/Metrics/cif_calibration_cause_1_NCC_ncontrols_100.png}
        \caption{Stratified CIF Calibration ($m=100$)}
    \end{subfigure}
    \hfill
    \begin{subfigure}[b]{0.48\textwidth}
        \includegraphics[width=\textwidth]{images/fh_results/NCC/Metrics/subject_curves_cause_1_cif_NCC_ncontrols_100.png}
        \caption{Marginal Calibration ($m=100$)}
    \end{subfigure}
    
    \vskip 1\baselineskip
    
    \begin{subfigure}[b]{0.48\textwidth}
        \includegraphics[width=\textwidth]{images/fh_results/NCC/Metrics/cif_calibration_cause_1_NCC_ncontrols_5.png}
        \caption{Stratified CIF Calibration ($m=5$)}
    \end{subfigure}
    \hfill
    \begin{subfigure}[b]{0.48\textwidth}
        \includegraphics[width=\textwidth]{images/fh_results/NCC/Metrics/subject_curves_cause_1_cif_NCC_ncontrols_5.png}
        \caption{Marginal Calibration ($m=5$)}
    \end{subfigure}
    \caption{Calibration diagnostics for NCC on the Framingham dataset. Both the optimal control pool ($m=100$, top) and the heavily constrained pool ($m=5$, bottom) maintain excellent absolute calibration. The high censoring rate inherently limits the variance of the background risk-set, preventing the IPW collapse observed in the Synthetic and PBC datasets.}
    \label{fig:fh_ncc_calibration}
\end{figure}

Both the optimal configuration ($m=100$) and the heavily constrained configuration ($m=5$) successfully track the empirical Aalen-Johansen baseline without experiencing probability collapse. The Stratified CIF curves remain highly calibrated, and the Marginal Calibration Gap is negligible for both settings. 

This resilience is a direct artifact of the data structure. Because 84.4\% of the Framingham cohort is censored (meaning they do not experience an event), the background risk-set is highly uniform and dense with "survivors". When the NCC sampler draws $m=5$ random controls, those controls are almost certainly censored patients representing the true low-hazard baseline. Therefore, the uniform random sampling of NCC accurately approximates the population density without introducing the severe variance spikes that destroyed calibration in the denser PBC and Synthetic datasets.


\subsubsection{Static Case-Cohort (SCC)}

The application of the Static Case-Cohort (SCC) sampler to the Framingham dataset highlights a fundamental difference between static and dynamic sampling strategies under conditions of severe right-censoring. 

When a robust background cohort is maintained ($\alpha=0.5$), the SCC architecture performs optimally. As shown in the top row of Figure \ref{fig:fh_scc_calibration}, this configuration successfully replicates the Full Cohort baseline. The Stratified CIF curves track the empirical reality cleanly, and the Marginal Calibration plot shows near-perfect agreement between the model average and the population baseline.

\begin{figure}[htbp]
    \centering
    % Optimal Compression (alpha = 0.5)
    \begin{subfigure}[b]{0.48\textwidth}
        \includegraphics[width=\textwidth]{images/fh_results/SCC/Metrics/cif_calibration_cause_1_SCC_subcohort_0.5.png}
        \caption{Stratified CIF Calibration ($\alpha=0.5$)}
    \end{subfigure}
    \hfill
    \begin{subfigure}[b]{0.48\textwidth}
        \includegraphics[width=\textwidth]{images/fh_results/SCC/Metrics/subject_curves_cause_1_cif_SCC_subcohort_0.5.png}
        \caption{Marginal Calibration ($\alpha=0.5$)}
    \end{subfigure}
    
    \vskip 1\baselineskip
    
    % High Compression (alpha = 0.05)
    \begin{subfigure}[b]{0.48\textwidth}
        \includegraphics[width=\textwidth]{images/fh_results/SCC/Metrics/cif_calibration_cause_1_SCC_subcohort_0.05.png}
        \caption{Stratified CIF Calibration ($\alpha=0.05$)}
    \end{subfigure}
    \hfill
    \begin{subfigure}[b]{0.48\textwidth}
        \includegraphics[width=\textwidth]{images/fh_results/SCC/Metrics/subject_curves_cause_1_cif_SCC_subcohort_0.05.png}
        \caption{Marginal Calibration ($\alpha=0.05$)}
    \end{subfigure}
    \caption{Calibration diagnostics for SCC on the Framingham dataset. While the optimal subcohort ($\alpha=0.5$) maintains baseline calibration, extreme compression ($\alpha=0.05$) reveals a structural vulnerability. The static subcohort fails to represent the temporal hazard over the long 3,500-day timeline, leading to an expanding calibration gap.}
    \label{fig:fh_scc_calibration}
\end{figure}

However, unlike the NCC sampler which proved highly resilient to extreme compression on this specific dataset the SCC sampler demonstrates a clear calibration breakdown when constrained to $\alpha=0.05$. In the bottom row of Figure \ref{fig:fh_scc_calibration}, the Stratified CIF curves for the higher risk strata begin to systematically underestimate the empirical Aalen-Johansen estimates. Consequently, the Marginal Calibration plot reveals a distinct and expanding calibration gap, as the mean model prediction detaches from the population baseline over time.

This divergence is mathematically illuminating. While the high censoring rate provided the NCC sampler with a dense pool of continuous "survivors" to sample from at each distinct event time, the SCC architecture is structurally rigid; it fixes its subcohort at $t=0$. In a long-horizon clinical study spanning nearly 10 years, drawing only a 5\% static sample at the baseline fails to provide sufficient covariate density to represent the evolving background hazard accurately. The IPW scaling becomes temporally skewed, proving that static sampling is more vulnerable to prolonged right-censoring than dynamic, time-matched sampling.

\subsubsection{Entropy Sampler}

The evaluation of the Entropy sampler on the heavily right-censored Framingham dataset exposes the complex interplay between dynamic risk-weighting and baseline hazard estimation. Because only 15.6\% of the Framingham cohort experiences the primary event, the background risk-set is overwhelmingly dominated by censored "survivors." 

When the temperature parameter is set to a highly greedy state ($\tau=0.5$), the Entropy algorithm aggressively prioritizes the rare, high-variance patients experiencing events, severely down-weighting the massive pool of censored observations. In a highly censored dataset, this over concentration proves detrimental to absolute calibration. 

\begin{figure}[htbp]
    \centering
    \begin{subfigure}[b]{0.48\textwidth}
        \includegraphics[width=\textwidth]{images/fh_results/Entropy/Metrics/cif_calibration_cause_1_Entropy_tau_0.5_compression_0.1.png}
        \caption{Stratified CIF Calibration ($\tau=0.5$)}
    \end{subfigure}
    \hfill
    \begin{subfigure}[b]{0.48\textwidth}
        \includegraphics[width=\textwidth]{images/fh_results/Entropy/Metrics/subject_curves_cause_1_cif_Entropy_tau_0.5_compression_0.1.png}
        \caption{Marginal Calibration ($\tau=0.5$)}
    \end{subfigure}
    \caption{Calibration diagnostics for a greedy Entropy sampler ($\tau=0.5$) on the Framingham dataset. By hyper-focusing on rare events and down-weighting censored observations, the IPW scaling becomes distorted. The model overestimates high-risk incidence (left, solid green line) and systematically underestimates the marginal population average (right).}
    \label{fig:fh_entropy_05_calibration}
\end{figure}

As illustrated in Figure \ref{fig:fh_entropy_05_calibration}, the $\tau=0.5$ configuration distorts the Inverse Probability Weighting (IPW) scaling. In the left panel, the model noticeably overestimates the true incidence for the high-risk stratum (solid green line exceeding the dashed empirical line). Concurrently, the Marginal Calibration diagnostic (right panel) reveals a distinct calibration gap, with the mean model prediction systematically underestimating the overall population baseline. By starving the neural network of the censored background population, the model fails to accurately anchor its absolute probability estimates.

To rectify this, the temperature parameter must be increased to soften the sampling distribution. Setting $\tau=5.0$ forces the Entropy sampler to behave less greedily, ensuring that while high-information events are still prioritized, a statistically viable volume of censored background patients is retained in every mini-batch. 

\begin{figure}[htbp]
    \centering
    \begin{subfigure}[b]{0.48\textwidth}
        \includegraphics[width=\textwidth]{images/fh_results/Entropy/Metrics/cif_calibration_cause_1_Entropy_tau_5_compression_0.1.png}
        \caption{Stratified CIF Calibration ($\tau=5.0$)}
    \end{subfigure}
    \hfill
    \begin{subfigure}[b]{0.48\textwidth}
        \includegraphics[width=\textwidth]{images/fh_results/Entropy/Metrics/subject_curves_cause_1_cif_Entropy_tau_5_compression_0.1.png}
        \caption{Marginal Calibration ($\tau=5.0$)}
    \end{subfigure}
    \caption{Calibration diagnostics for the optimally tuned Entropy sampler ($\tau=5.0$). Softening the probability distribution allows the model to retain sufficient censored observations to anchor the baseline hazard. The IPW mechanism is restored, yielding near-perfect Stratified CIF tracking and eliminating the Marginal Calibration Gap.}
    \label{fig:fh_entropy_5_calibration}
\end{figure}

Figure \ref{fig:fh_entropy_5_calibration} demonstrates the complete restoration of absolute calibration at $\tau=5.0$. The Stratified CIF curves now align perfectly with the empirical Aalen-Johansen estimates across all strata, eliminating the high-risk overestimation observed at $\tau=0.5$. Furthermore, the Marginal Calibration Gap is entirely closed, matching the structural integrity of the Full Cohort baseline.

This final evaluation solidifies the overarching operational framework for the CRISPNAM architecture. Across both small, noisy clinical trials (PBC) and heavily censored longitudinal studies (Framingham), highly greedy entropy sampling ($\tau=0.5$) corrupts IPW calibration. However, by treating the temperature parameter $\tau$ as an environmental hyperparameter tuning it upward ($\tau \in [1.0, 5.0]$) to flatten the sampling distribution in the presence of severe data sparsity or noise the Entropy sampler successfully matches Full Cohort predictive power and absolute calibration at a fraction of the computational cost.


\section{Feature Importance and Interpretability}


\section{Discussion}

also mention that our c-index and brier are close to the papers we used for refrnce 
\subsection{Computational Efficiency and Vectorization}

\chapter{Conclusion and Future Work}
\section{Summary of Findings}
\section{Limitations and Future Directions}

\bibliography{references}\addcontentsline{toc}{chapter}{References}

% ======================================================================
% APPENDICES & REFERENCES
% ======================================================================
\appendix
\chapter{Code Listings}
% Insert vectorization or sampler code here

\chapter{Supplementary Figures}
\section{Hardware Resource Utilization}

\ifdefined\includeHardwareResourceFigures


\begin{figure}[htbp]
    \centering
    % Row 1: Full Cohort
    \begin{subfigure}[b]{0.48\textwidth}
        \includegraphics[width=\textwidth]{images/Synth_results/FULL/fold_1_resource_usage.png}
        \caption{Full Cohort: CPU}
    \end{subfigure}
    \hfill
    \begin{subfigure}[b]{0.48\textwidth}
        \includegraphics[width=\textwidth]{images/Synth_results/FULL/fold_1_memory_usage.png}
        \caption{Full Cohort: Memory}
    \end{subfigure}
    \vskip\baselineskip
    
    % Row 2: NCC
    \begin{subfigure}[b]{0.48\textwidth}
        \includegraphics[width=\textwidth]{images/Synth_results/NCC/fold_1_resource_usage_NCC_ncontrols_5.png}
        \caption{NCC ($m=5$): CPU}
    \end{subfigure}
    \hfill
    \begin{subfigure}[b]{0.48\textwidth}
        \includegraphics[width=\textwidth]{images/Synth_results/NCC/fold_1_memory_usage_NCC_ncontrols_5.png}
        \caption{NCC ($m=5$): Memory}
    \end{subfigure}
    \vskip\baselineskip
    
    % Row 3: SCC
    \begin{subfigure}[b]{0.48\textwidth}
        \includegraphics[width=\textwidth]{images/Synth_results/SCC/fold_1_resource_usage_SCC_subcohort_0.5.png}
        \caption{SCC ($\alpha=0.5$): CPU}
    \end{subfigure}
    \hfill
    \begin{subfigure}[b]{0.48\textwidth}
        \includegraphics[width=\textwidth]{images/Synth_results/SCC/fold_1_memory_usage_SCC_subcohort_0.5.png}
        \caption{SCC ($\alpha=0.5$): Memory}
    \end{subfigure}
    \vskip\baselineskip

    % Row 4: Entropy
    \begin{subfigure}[b]{0.48\textwidth}
        \includegraphics[width=\textwidth]{images/Synth_results/ENTROPY/fold_1_resource_usage_Entropy_tau_0.5_compression_0.1.png}
        \caption{Entropy ($\tau=0.5$): CPU}
    \end{subfigure}
    \hfill
    \begin{subfigure}[b]{0.48\textwidth}
        \includegraphics[width=\textwidth]{images/Synth_results/ENTROPY/fold_1_memory_usage_Entropy_tau_0.5_compression_0.1.png}
        \caption{Entropy ($\tau=0.5$): Memory}
    \end{subfigure}
    
    \caption{Hardware resource utilization (CPU and Memory) over time for the synthetic dataset ($N=15,000$). While sampling strategies introduce a slight memory overhead during initialization, they significantly compress the execution time, with Entropy sampling achieving a >50\% reduction in training duration compared to the full cohort.}
    \label{fig:resource_profiling}
\end{figure}
\fi

\end{document}